\documentclass[12pt]{article}

\usepackage{verbatim,color,amssymb}
\usepackage{amsmath}					
\usepackage{amsthm}		
\usepackage{bm}
\usepackage{dsfont}			
\usepackage{natbib}
\usepackage{multirow}
\usepackage{setspace}
\usepackage[mathscr]{euscript}
\usepackage{fancyhdr}
\usepackage{enumitem}
\usepackage{graphicx}
\usepackage{geometry}
\usepackage{xcolor}
\usepackage{hyperref}
\usepackage{bbm}

%\usepackage[bookmarks=false]{hyperref}

%\usepackage{multibib}
%\newcites{latex}{References}

\usepackage{tabularx, booktabs}
\newcolumntype{Y}{>{\centering\arraybackslash}X}
\usepackage{lscape}


\usepackage{subfig}
%\usepackage{caption}
%\captionsetup[subfigure]{labelformat=parens,labelsep=space,font=small}
\usepackage{float}
\usepackage{lineno}


\def\pdfshellescape{1}

\setlength{\textheight}{9in}
\setlength{\textwidth}{6in}
\setlength{\topmargin}{-36pt}
\setlength{\oddsidemargin}{15pt}
\setlength{\evensidemargin}{0pt}
\tolerance=500
\renewcommand{\baselinestretch}{1.5}

% Gio comment
\newcommand{\gio}{\color{red}} 
\newcommand{\ech}{\color{black}}

\newcommand*{\swap}[2]{#2#1} %% Swap order (used to change order referencees)

%%%%%%%%%%%%%%%
% Begin New Definitions  %%
%%%%%%%%%%%%%%%
%\iffalse
\newtheorem{Thm}{\underline{\bf Theorem}}
\newtheorem{Assmp}{{\bf Assumption}}
\newtheorem{Cond}{\underline{\bf Conditions}}
\newtheorem{Proof}{Proof}
\newtheorem*{Proof*}{Proof}
\newtheorem{Mth}{Main Theorem}
\newtheorem{Def}{Definition}
\newtheorem{Rem}{\underline{\bf Remark}}
\newtheorem{Qes}{Question}
\newtheorem{Prop}{\underline{\bf Proposition}}
\newtheorem{Lem}{\underline{\bf Lemma}}
\newtheorem{Cor}{\underline{\bf Corollary}}
\newtheorem{Exa}{Example}
\newtheorem{Eq}{Equation}
%\fi

\newcommand{\MyProof}{\noindent\textbf{Proof. }}
\newcommand{\mytilde}[2][1.2]{
	\mathrel{\overset{\mbox{\tiny #2}}{\scalebox{#1}[1]{$\sim$}}}}

\def\nN{\mathbb{N}}
\def\rR{\mathbb{R}}
\def \R {\mathbb{R} }
\def\eE{\mathbb{E}}
\def \d {\mathrm{d}}

\def\L{{\cal L}}
\def\B{{\cal B}}
\def\C{{\cal C}}
\def\D{{\cal D}}
\def\E{{\cal E}}
\def\F{{\cal F}}
\def\G{{\cal G}}
\def\H{{\cal H}}
\def\I{{\cal I}}
\def\K{{\cal K}}
\def\M{{\cal M}}
\def\N{{\cal N}}
\def\calP{{\cal P}}
\def\S{{\cal S}}
\def\T{{\cal T}}
\def\U{{\cal U}}
\def\W{{\cal W}}
\def\V{{\cal V}}
\def\X{{\cal X}}
\def\Z{{\cal Z}}
\def\Y{{\cal Y}}
\def\sumi{\sum_{i=1}^n}

\def\scrC{{\mathscr{C}}}


\def\diag{\hbox{diag}}
\def\Ind{\hbox{I}}
\def\wh{\widehat}
\def\wt{\widetilde}
\def\wb{\breve}
\def\AIC{\hbox{AIC}}
\def\BIC{\hbox{BIC}}
\def\diag{\hbox{diag}}
\def\log{\hbox{log}}
\def\bias{\hbox{bias}}
\def\Siuu{\boldSigma_{i,uu}}
\def\whT{\widehat{\Theta}}
\def\var{\hbox{var}}
\def\cov{\hbox{cov}}
\def\corr{\hbox{corr}}
\def\sign{\hbox{sign}}
\def\trace{\hbox{trace}}
\def\naive{\hbox{naive}}
\def\vect{\hbox{vec}}


\def\Bern{\hbox{Bernoulli}}
\def\Beta{\hbox{Beta}}
\def\DE{\hbox{DE}}
\def\Dir{\hbox{Dir}}
\def\DP{\hbox{DP}}
\def\Exp{\hbox{Exp}}
\def\gIGs{\hbox{g-Inv-Gs}}
\def\Ga{\hbox{Ga}}
\def\GP{\hbox{GP}}
\def\GEM{\hbox{GEM}}
\def\HC{\hbox{C}^{+}}
\def\IGs{\hbox{Inv-Gs}}
\def\IG{\hbox{Inv-Ga}}
\def\IW{\hbox{IW}}
\def\MVN{\hbox{MVN}}
\def\MVL{\hbox{MVL}}
\def\MVT{\hbox{MVT}}
\def\Normal{\hbox{Normal}}
\def\FN{\hbox{Folded-Normal}}
\def\HN{\hbox{Half-Normal}}
\def\LN{\hbox{Log-Normal}}
\def\Poi{\hbox{Poi}}
\def\SB{\hbox{SB}}
\def\TN{\hbox{TN}}
\def\Unif{\hbox{Unif}}
\def\Mult{\hbox{Mult}}
\def\Wish{\hbox{W}}
\def\HOHMM{\hbox{HOHMM}}

\def\ANNALS{{\it Annals of Statistics}}
\def\ANNALSP{{\it Annals of Probability}}
\def\ANNALSMS{{\it Annals of Mathematical Statistics}}
\def\ANNALSAS{{\it Annals of Applied Statistics}}
\def\ANNALSISM{{\it Annals of the Institute of Statistical Mathematics}}
\def\AJE{{\it American Journal of Epidemiology}}
\def\ANIPS{{\it Advances in Neural Information Processing Systems}}
\def\APLS{{\it Applied Statistics}}
\def\BA{{\it Bayesian Analysis}}
\def\BRNL{{\it Bernoulli}}
\def\BIOK{{\it Biometrika}}
\def\BIOS{{\it Biostatistics}}
\def\BMCS{{\it Biometrics}}
\def\BMCMIDM{{\it BMC Medical Informatics and decision making}}
\def\BIOINF{{\it Bioinformatics}}
\def\CANADAJS{{\it Canadian Journal of Statistics}}
\def\CG{{\it Current Genomics}}
\def\CDA{{\it Computational Statistics \& Data Analysis}}
\def\COMMS{{\it Communications in Statistics, Series A}}
\def\COMMS{{\it Communications in Statistics, Theory \& Methods}}
\def\COMMSS{{\it Communications in Statistics - Simulation}}
\def\COMMSSC{{\it Communications in Statistics - Simulation and Computation}}
\def\EJS{{\it Electronic Journal of Statistics}}
\def\ECMK{{\it Econometrica}}
\def\ECTH{{\it Econometric Theory}}
\def\GENEP{{\it Genetic Epidemiology}}
\def\JASA{{\it Journal of the American Statistical Association}}
\def\JRSSB{{\it Journal of the Royal Statistical Society, Series B}}
\def\JRSSC{{\it Journal of the Royal Statistical Society, Series C}}
\def\JQT{{\it Journal of Quality Technology}}
\def\JCGS{{\it Journal of Computational and Graphical Statistics}}
\def\JCB{{\it Journal of Computational Biology}}
\def\JAMA{{\it Journal of the American Medical Association}}
\def\JNUTR{{\it Journal of Nutrition}}
\def\JABES{{\it Journal of Agricultural, Biological and Environmental Statistics}}
\def\JBES{{\it Journal of Business and Economic Statistics}}
\def\JSPI{{\it Journal of Statistical Planning \& Inference}}
\def\JMA{{\it Journal of Multivariate Analysis}}
\def\JNS{{\it Journal of Nonparametric Statistics}}
\def\JSS{{\it Journal of Statistical Software}}
\def\JECM{{\it Journal of Econometrics}}
\def\IEEE{{\it IEEE}}
\def\IEEESPL{{\it IEEE Signal Processing Letters}}
\def\IEEETIT{{\it IEEE Transactions on Information Theory}}
\def\LETTERS{{\it Letters in Probability and Statistics}}
\def\P_25_ICML{{\it Proceedings of the 25th international conference on Machine learning}}
\def\PLoSCB{{\it PloS Computational Biology}}
\def\STIM{{\it Statistics in Medicine}}
\def\SCAN{{\it Scandinavian Journal of Statistics}}
\def\SMMR{{\it Statistical Methods in Medical Research}}
\def\SNKH{{\it Sankhy\={a}: The Indian Journal of Statistics}}
\def\STIM{{\it Statistics in Medicine}}
\def\STATMED{{\it Statistics in Medicine}}
\def\STATSCI{{\it Statistical Science}}
\def\SSNC{{\it Statistica Sinica}}
\def\SaC{{\it Statistics and Computing}}
\def\STATSCI{{\it Statistical Science}}
\def\TECH{{\it Technometrics}}


\def\VS{{\vskip 3mm\noindent}}
\def\refhg{\hangindent=20pt\hangafter=1}
\def\refmark{\par\vskip 2mm\noindent\refhg}
\def\itemitem{\par\indent \hangindent2\parindent \textindent}
\def\refhg{\hangindent=20pt\hangafter=1}
\def\refmark{\par\vskip 2mm\noindent\refhg}
\def\povr{\buildrel p\over\longrightarrow}
\def\ccdot{{\bullet}}
\def\bse{\begin{eqnarray*}}
	\def\ese{\end{eqnarray*}}
\def\be{\begin{eqnarray}}
	\def\ee{\end{eqnarray}}
\def\bq{\begin{equation}}
	\def\eq{\end{equation}}
\def\pr{\hbox{pr}}
\def\wh{\widehat}


\def\trans{^{\rm T}}
\def\transs{\hbox{\footnotesize T}}
\def\rank{{\hbox{rank}}}
\def\myalpha{{\cal A}}
\def\th{^{th}}
\def\bone{{\mathbf 1}}
\def\Data{{\hbox{Data}}}


\def\ba{{\mathbf a}}
\def\bA{{\mathbf A}}
\def\bb{{\mathbf b}}
\def\bB{{\mathbf B}}
\def\bc{{\mathbf c}}
\def\bC{{\mathbf C}}
\def\bd{{\mathbf d}}
\def\bD{{\mathbf D}}
\def\b1e{{\mathbf e}}
\def\b1f{{\mathbf f}}
\def\bF{{\mathbf F}}
\def\bG{{\mathbf G}}
\def\bg{{\mathbf g}}
\def\bh{{\mathbf h}}
\def\bH{{\mathbf H}}
\def\bI{{\mathbf I}}
\def\bk{{\mathbf k}}
\def\bM{{\mathbf M}}
\def\bn{{\mathbf n}}
\def\bN{{\mathbf N}}
\def\bp{{\mathbf p}}
\def\bP{{\mathbf P}}
\def\br{{\mathbf r}}
\def\bR{{\mathbf R}}
\def\bs{{\mathbf s}}
\def\bS{{\mathbf S}}
\def\bt{{\mathbf t}}
\def\bu{{\mathbf u}}
\def\bU{{\mathbf U}}
\def\bv{{\mathbf v}}
\def\bV{{\mathbf V}}
\def\bw{{\mathbf w}}
\def\bW{{\mathbf W}}
\def\bx{{\mathbf x}}
\def\bX{{\mathbf X}}
\def\by{{\mathbf y}}
\def\bY{{\mathbf Y}}
\def\bz{{\mathbf z}}
\def\bZ{{\mathbf Z}}
\def\bS{{\mathbf S}}
\def\bzero{{\mathbf 0}}

\def\whT{\widehat{\Theta}}
\def\te{\widetilde{e}}
\def\te{\widetilde{\epsilon}}
\def\tp{\widetilde{p}}
\def\tv{\widetilde{v}}
\def\tmu{\widetilde{\mu}}
\def\tsigma{\widetilde{\sigma}}

\newcommand{\etal}{\emph{et al.~}}
\newcommand{\etam}{\mbox{\boldmath $\eta$}}
\newcommand{\bmu}{\mbox{\boldmath $\mu$}}
\newcommand{\bDelta}{\mbox{\boldmath $\Delta$}}
\newcommand{\bphi}{\mbox{\boldmath $\phi$}}
\newcommand{\bpi}{\mbox{\boldmath $\pi$}}
\newcommand{\bPi}{\mbox{\boldmath $\Pi$}}
\newcommand{\bxi}{\mbox{\boldmath $\xi$}}
\newcommand{\bepsilon}{\mbox{\boldmath $\epsilon$}}
\newcommand{\btheta}{\mbox{\boldmath $\theta$}}
\newcommand{\bTheta}{\mbox{\boldmath $\Theta$}}
\newcommand{\bbeta}{\mbox{\boldmath $\beta$}}
\newcommand{\bgamma}{\mbox{\boldmath $\gamma$}}
\newcommand{\bzeta}{\mbox{\boldmath $\zeta$}}
\newcommand{\bsigma}{\mbox{\boldmath $\sigma$}}
\newcommand{\bSigma}{\mbox{\boldmath $\Sigma$}}
\newcommand{\balpha}{\mbox{\boldmath $\alpha$}}
\newcommand{\bomega}{\mbox{\boldmath $\omega$}}
\newcommand{\blambda}{\mbox{\boldmath $\lambda$}}
\newcommand{\bLambda}{\mbox{\boldmath $\Lambda$}}
\newcommand{\bOmega}{\mbox{\boldmath $\Omega$}}
\newcommand{\bPsi}{\mbox{\boldmath $\Psi$}}
\newcommand{\bpsi}{\mbox{\boldmath $\psi$}}
\newcommand{\bGamma}{\mbox{\boldmath $\Gamma$}}
\newcommand{\btau}{\mbox{\boldmath $\tau$}}

\newcommand{\abs}[1]{\left\vert#1\right\vert}
\newcommand{\norm}[1]{\left\Vert#1\right\Vert}
\newcommand*{\hatequal}{\mathrel{\widehat=}}


\def\simind{\stackrel{\mbox{\scriptsize{ind}}}{\sim}}
\def\simiid{\stackrel{\mbox{\scriptsize{iid}}}{\sim}}
\def\eqd{\stackrel{\mbox{\scriptsize{d}}}{=}}
\def\as{\stackrel{\mbox{\scriptsize{a.s.}}}{\longrightarrow}}
\def\qed{\hfill $\Box$}

\renewcommand\footnoterule{\kern-3pt \hrule \textwidth 2in \kern 2.6pt}

\newcommand*{\myalign}[2]{\multicolumn{1}{#1}{#2}}

\newcommand\Algphase[1]{%
	\vspace*{-.5\baselineskip}\Statex\hspace*{\dimexpr-\algorithmicindent-2pt\relax}\rule{\textwidth}{0.4pt}%
	\Statex\hspace*{-\algorithmicindent}\textbf{#1}%
	\vspace*{-.5\baselineskip}\Statex\hspace*{\dimexpr-\algorithmicindent-2pt\relax}\rule{\textwidth}{0.4pt}%
}

\def\colblue#1{\textcolor{blue}{\bf #1}}
\def\colred#1{\textcolor{red}{ #1}}
\def\colredb#1{\textcolor{red}{\bf #1}}
\def\highlight#1{\underline{\textcolor{red}{#1}}}
\def\boxit#1{\vbox{\hrule\hbox{\vrule\kern6pt \vbox{\kern6pt \textcolor{blue}{#1}\kern6pt}\kern6pt\vrule}\hrule}}
\def\grcomment#1{\vskip 2mm\boxit{\vskip 2mm{\color{blue}\bf#1} {\color{black}\bf -- Gio\vskip 2mm}}\vskip 2mm}

\def\authorfootnote#1{{\let\thefootnote\relax\footnotetext{#1}}}



%%%%%%%%%%%%%%
% End New Definitions  %%
%%%%%%%%%%%%%%

\allowdisplaybreaks


\pdfminorversion=4

\begin{document}\newpage


\rhead{\bfseries\thepage}
\lhead{\bfseries SUPPLEMENTARY MATERIALS}

\baselineskip 25pt
\begin{center}
{\LARGE{Supplementary Materials for\\} 
\bf  Clustering consistency with\\ Dirichlet process mixtures
}
\end{center}


\setcounter{equation}{0}
\setcounter{page}{1}
\setcounter{table}{1}
\setcounter{figure}{0}
\setcounter{Thm}{0}
\setcounter{Lem}{0}
\setcounter{Cor}{0}
\setcounter{Prop}{0}
\setcounter{section}{0}
\numberwithin{table}{section}
\renewcommand{\theequation}{S.\arabic{equation}}
\renewcommand{\thesubsection}{S.\arabic{section}.\arabic{subsection}}
\renewcommand{\theThm}{S.\arabic{Thm}}
\renewcommand{\theCor}{S.\arabic{Cor}}
\renewcommand{\theProp}{S.\arabic{Prop}}
\renewcommand{\theLem}{S.\arabic{Lem}}
\renewcommand{\thesection}{S.\arabic{section}}
\renewcommand{\thepage}{S.\arabic{page}}
\renewcommand{\thetable}{S.\arabic{table}}
\renewcommand{\thefigure}{S.\arabic{figure}}

\vspace{0cm}

%\iffalse
	\small
\baselineskip=14pt
	\begin{center}
		Filippo Ascolani$^{a,b}$ (filippo.ascolani@phd.unibocconi.it)\\
		Antonio Lijoi$^{a,b}$ (antonio.lijoi@unibocconi.it)\\
		Giovanni Rebaudo$^{b,c}$ (giovanni.rebaudo@austin.utexas.edu)\\
		Giacomo Zanella$^{a,b}$ (giacomo.zanella@unibocconi.it)
	
	%%%%%%%%%%%%%%%%%%%%%%%%%%%%%
	\vskip 3mm
	$^{a}$Department of Decision Sciences, Bocconi University,\\
	via R\"oentgen 1, 20136 Milan, Italy
	\vskip 4pt 
	$^{b}$Bocconi Institute for Data Science and Analytics, Bocconi University,\\ 
	Via Röntgen 1, 20136 Milan, Italy
	\vskip 4pt 
	$^{c}$Department of Statistics and Data Sciences,
	University of Texas at Austin,\\
	105 East 24th Street D9800, Austin, TX 78712, USA\\	
\end{center}

\vskip 10mm
\baselineskip=\gnat

\section*{Summary}
This document contains the proofs of all the results in the main document. In order to follow the logical lines of the arguments, we present the proofs of Sections \ref{sec:Methodology} and \ref{sec:consistency_results}, in this order.

\section{Proof of Lemma \ref{sufficient}}
\begin{proof}
The result immediately follows upon noting that
\[
\text{pr}(K_n = t \mid X_{1:n})=  \left\{1+\sum_{s \neq t} \frac{\text{pr}(K_n = s \mid X_{1:n} )}{\text{pr}(K_n = t \mid X_{1:n})}\right\}^{-1}.
\]
\end{proof}



\section{Proof of Proposition \ref{prop:Prior}}
By assumptions $A1$ and $A2$ there exist $\epsilon, \delta, \beta > 0$ such that
\begin{equation}\label{motivation}
\frac{1}{\delta^2}\frac{\int_0^\epsilon \frac{\alpha^{t+s+\beta}}{\alpha^{(n)}}\, \mathrm{d}\alpha}{\int_0^\epsilon \frac{\alpha^{t+\beta}}{\alpha^{(n)}} \, \mathrm{d}\alpha}  \leq
\frac{\int_0^\epsilon \frac{\alpha^{t+s}}{\alpha^{(n)}}\pi(\alpha)\, \mathrm{d}\alpha}{\int_0^\epsilon \frac{\alpha^t}{\alpha^{(n)}}\pi(\alpha) \, \mathrm{d}\alpha}  \leq \delta^2\frac{\int_0^\epsilon \frac{\alpha^{t+s+\beta}}{\alpha^{(n)}}\, \mathrm{d}\alpha}{\int_0^\epsilon \frac{\alpha^{t+\beta}}{\alpha^{(n)}} \, \mathrm{d}\alpha}\,.
\end{equation}
Notice that, if assumption $A2$ holds for $\epsilon \geq 1$, it holds also for $\epsilon < 1$. Thus, without loss of generality, we will assume $\epsilon < 1$ and the main object of interest will be
\[
%\frac{\int_0^\epsilon \frac{\alpha^{s+1}}{\alpha^{(n)}} \, \pi(\alpha)\mathrm{d}\alpha}{\int_0^\epsilon \frac{\alpha}{\alpha^{(n)}} \, \pi(\alpha)\mathrm{d}\alpha} 
%%= 
%%\int_0^\epsilon \frac{ \frac{\alpha^{s+1}}{\alpha^{(n)}} \, \pi(\alpha)\mathrm{d}\alpha}{\int_0^\epsilon \frac{x}{x^{(n)}}\pi(x) \, \mathrm{d}x} 
%=
E_{n}(\alpha^{s})= \int_0^\epsilon \alpha^{s} p_{n}(\alpha) \mathrm{d}\alpha,
\]
where $E_n$ denotes the expected value with respect to the probability distribution with density
\begin{equation}\label{n_density}
p_n(\alpha) = \frac{f_n(\alpha)}{\int_0^\epsilon f_n(x)\, \mathrm{d}x}, \quad 
f_n(x) = \frac{x^{t+\beta}}{x^{(n)}}\:\mathbbm{1}_{(0,\epsilon)}(x),
\end{equation}
where $\mathbbm{1}_A$ stands for the indicator function of set $A$. We now provide three lemmas that will be useful to prove Proposition 1.

\begin{Lem}\label{lemma:stoch_ord}
Let $f$ and $g$ be two pdf's on $\R$ such that $g(x)/f(x)$ is non-decreasing in $x$.
Then $\int h(x)f(x)\mathrm{d}x\leq \int h(x)g(x)\mathrm{d}x$ for every non-decreasing $h:\R\to\R$.
\end{Lem}
\begin{proof}
Let $X\sim f$ and $Y\sim g$.
Since $g(x)/f(x)$ is non-decreasing we have $g(x_0)f(x_1) \le g(x_1)f(x_0)$ for every $x_0 < x_1$. 
Thus we have %by monotonicity of integral we get
\begin{align*}
F_{Y}(x_1) f(x_1) &= \int_{-\infty}^{x_1} g(x_0)f(x_1) \mathrm{d}x_0 \le  \int_{-\infty}^{x_1} g(x_1)f(x_0) \mathrm{d}x_0 = F_{X}(x_1) g(x_1)
\end{align*}
and
\begin{align*}
%F_{Y}(x_1) f(x_1) &= \int_{-\infty}^{x_1} g(x_0)f(x_1) \mathrm{d}x_0 \le  \int_{-\infty}^{x_1} g(x_1)f(x_0) \mathrm{d}x_0 = F_{X}(x_1) g(x_1),\\
\{1-F_{X}(x_0)\} g(x_0) &= \int_{x_0}^{\infty} g(x_0)f(x_1) \mathrm{d}x_1 \le  \int_{x_0}^{\infty} g(x_1)f(x_0) \mathrm{d}x_1 = \{1-F_{Y}(x_0) \} f(x_0).
\end{align*}
It follows
%\begin{align*}
%    \frac{F_{Y}(x)}{F_{X}(x)} \le \frac{g(x)}{f(x)}, \quad \frac{1-F_{Y}(x)}{1-F_{X}(x)} \textcolor{red}{\ge} \frac{g(x)}{f(x)}.
%\end{align*}
\begin{align*}
\frac{F_{Y}(x)}{F_{X}(x)} \le \frac{g(x)}{f(x)}\le\frac{1-F_{Y}(x)}{1-F_{X}(x)}\,,% \textcolor{red}{\ge} \frac{g(x)}{f(x)}.
\end{align*}
for every $x\in \R$,
which implies
\begin{align*}
\frac{F_{Y}(x)}{1-F_{Y}(x)} \le \frac{F_{X}(x)}{1-F_{X}(x)}\,.% \textcolor{red}{\ge} \frac{g(x)}{f(x)}.
\end{align*}
%The last two inequalities imply that $F_Y(x)\{1-F_X(x)\} \le F_X(x) \{1-F_Y(x) \}$. 
Thus, $Y$ stochastically dominates $X$, i.e.\ the corresponding cdf's satisfy $F_Y(x)~\leq~F_X(x)$ for every $x\in \R$, which implies that $E\{h(X)\}\leq E\{h(Y)\}$ for every non-decreasing $h$.
\end{proof}

\begin{Lem}\label{main}
Under assumptions $A1$ and $A2$, for every $n-t > s \geq 1$ it holds
\[
\frac{\gamma[t+s+\beta,\epsilon\{\log(n)+1\}]}{\delta^2\gamma[t+\beta,\epsilon\{\log(n)+1\}]}\{\log(n)+1\}^{-s}
\leq
\frac{\int_0^\epsilon \frac{\alpha^{t+s}}{\alpha^{(n)}}\pi(\alpha)\, \mathrm{d}\alpha}{\int_0^\epsilon \frac{\alpha^t}{\alpha^{(n)}}\pi(\alpha) \, \mathrm{d}\alpha} 
\leq
\frac{\delta^2 \gamma\{t+s+\beta,\epsilon\log(n)\}}{\gamma\{t+\beta,\epsilon\log(n)\}}\{\log(n)/(1+\epsilon)\}^{-s},
\]
where $\gamma(x,y)$ is the lower incomplete Gamma function and we recall that $\epsilon, \delta, \beta >0$ are such that for every $ \alpha \in (0, \epsilon)$ it holds $\frac{1}{\delta}\alpha^\beta \leq \pi(\alpha) \leq \delta \alpha^\beta$.
\end{Lem}
\begin{proof}
By \eqref{motivation} it suffices to find suitable bounds of $E_n(\alpha^{s})$.
For the upper inequality we apply Lemma \ref{lemma:stoch_ord} with $f=p_n$, $g(\alpha)\propto (cn)^{-\alpha}\alpha^{t+\beta-1}\mathbbm{1}_{\left(\alpha\in[0,\epsilon]\right)}$ with $c=(1+\epsilon)^{-1}$ and $h(\alpha)=\alpha^s$.
To verify that $g(\alpha)/p_n(\alpha)$ is non-decreasing for $\alpha\in(0,\epsilon]$ we compute 
\begin{align*}
\frac{\mathrm{d}}{\mathrm{d}\alpha}\log\left\{\frac{g(\alpha)}{p_n(\alpha)}\right\}=&
-\log \left(\frac{n}{1+\epsilon}\right)+\sum_{i=1}^{n-1}\frac{1}{\alpha+i}\\
\geq&
-\log \left(\frac{n+\epsilon}{1+\epsilon}\right)+\sum_{i=1}^{n-1}\frac{1}{i+\epsilon}
\geq 0,
\end{align*}
where the last inequality follows from 
\[
\int_1^{k}\frac{1}{x+\epsilon} \, dx < \sum_{i = 1}^{k-1}\frac{1}{i+\epsilon}
\]
for every $k > 1$. Thus, since $h(\alpha)=\alpha^s$ is non-decreasing in $\alpha$ it follows by Lemma \ref{lemma:stoch_ord} that
$$
\begin{aligned}
E_n(\alpha^{s})
\leq&
\frac{\int_0^\epsilon \alpha^{t+s+\beta-1}(cn)^{-\alpha} \mathrm{d}\alpha}{\int_0^\epsilon \alpha^{t+\beta-1}(cn)^{-\alpha}\, \mathrm{d}\alpha} 
=
\frac{\{\log(cn)\}^{-s}\int_0^{\epsilon\log(cn)} z^{t+s+\beta-1}e^{-z} \mathrm{d}z}{\int_0^{\epsilon\log(cn)} z^{t+\beta-1}e^{-z}\, \mathrm{d}z} 
\\
=&
\frac{\{\log(cn)\}^{-s}\gamma\{t+s+\beta,\epsilon\log(cn)\}}{\gamma\{t+\beta,\epsilon\log(cn)\}}.
\end{aligned}
$$
The lower bound again follows from Lemma \ref{lemma:stoch_ord} with $f(\alpha) \propto (e n)^{-\alpha}\alpha^{t+\beta-1}\mathbbm{1}_{\left(\alpha\in[0,\epsilon]\right)}$, $g(\alpha)=p_n(\alpha)$ and $h(\alpha)=\alpha^s$.
To verify that $p_n(\alpha)/f(\alpha)$ is non-decreasing for $\alpha\in(0,\epsilon]$ we compute
\begin{align*}
\frac{\mathrm{d}}{\mathrm{d}\alpha}\log\left\{\frac{p_n(\alpha)}{f(\alpha)}\right\}=&
-\sum_{i=1}^{n-1}\frac{1}{\alpha+i}+\log (n)+1\\
\geq&
-\sum_{i=1}^{n-1}\frac{1}{i}+\log (n)+1
\geq 0,
\end{align*}
where the last inequality follows from
\[
\sum_{i = 1}^{k}\frac{1}{i} \leq \log(k)+1
\]
for every $k \geq 1$. Thus, since $h(\alpha)=\alpha^s$ is non-decreasing in $\alpha$, we have 
$$
\begin{aligned}
E_n(\alpha^{s})
\geq&
\frac{\int_0^\epsilon \alpha^{t+s+\beta-1}(e n)^{-\alpha} \mathrm{d}\alpha}{\int_0^\epsilon \alpha^{t+\beta-1}(e n)^{-\alpha}\, \mathrm{d}\alpha} 
=
\frac{\{\log(en)\}^{-s}\int_0^{\epsilon\log(en)} z^{t+s+\beta-1}e^{-z} \mathrm{d}z}{\int_0^{\epsilon\log(en)} z^{t+\beta-1}e^{-z}\, \mathrm{d}z} \\
&=
\frac{\{\log(en)\}^{-s}\gamma\{t+s+\beta,\epsilon\log(en)\}}{\gamma\{t+\beta,\epsilon\log(en)\}}.
\end{aligned}
$$
The proof is completed by combining the bounds with \eqref{motivation}.
\end{proof}

\begin{Lem}\label{lower_tech}
For every $\epsilon >0$, there exists $M > 0$ such that, for every $n \geq 1$, it holds
\[
M \,\int_0^\epsilon \frac{\alpha^t}{\alpha^{(n)}}\, \pi(\alpha) \, \mathrm{d}\alpha \geq \int_\epsilon^\infty \frac{\alpha^t}{\alpha^{(n)}}\, \pi(\alpha) \, \mathrm{d}\alpha\,.
\]
\end{Lem}
\begin{proof}
Define $p = \frac{\int_\epsilon^\infty \alpha^t \pi(\alpha) \,\mathrm{d}\alpha}{\int_0^{\frac{\epsilon}{2}} \alpha^t \pi(\alpha) \,\mathrm{d}\alpha}$. Then
\begin{align*}
\int_0^\epsilon \frac{\alpha^t}{\alpha^{(n)}}\, \pi(\alpha) \, \mathrm{d}\alpha - \int_\epsilon^\infty \frac{\alpha^t}{\alpha^{(n)}}\, \pi(\alpha) \, \mathrm{d}\alpha=&
\int_0^{\epsilon} \frac{\alpha^t}{\alpha^{(n)}}\, \pi(\alpha) \, \mathrm{d}\alpha - \int_0^{\frac{\epsilon}{2}} p\frac{\alpha^t}{\epsilon^{(n)}}\, \pi(\alpha) \, \mathrm{d}\alpha\\
\geq& \int_0^{\frac{\epsilon}{2}} \frac{\alpha^t}{\alpha^{(n)}}\, \pi(\alpha) \, \mathrm{d}\alpha - \int_0^{\frac{\epsilon}{2}} p\frac{\alpha^t}{\epsilon^{(n)}}\, \pi(\alpha) \, \mathrm{d}\alpha.
\end{align*}
Choose $m$ such that $\left(\frac{\epsilon}{2}\right)^{(m)} < \frac{\epsilon^{(m)}}{p}$, which is always possible because $\left\{\epsilon^{(m)}\right\}^{-1} \left(\frac{\epsilon}{2}\right)^{(m)}\to 0$ as $m\to\infty$. Thus
\[
\int_0^\epsilon \frac{\alpha^t}{\alpha^{(n)}}\, \pi(\alpha) \, \mathrm{d}\alpha \geq \int_\epsilon^\infty \frac{\alpha^t}{\alpha^{(n)}}\, \pi(\alpha) \, \mathrm{d}\alpha, \quad n \geq m
\]
and it suffices to set $M = \max (P,1)$ with
\[
P = \max_{1 \leq i \leq m} \left\{ \frac{\int_\epsilon^\infty \frac{\alpha^t}{\alpha^{(i)}}\, \pi(\alpha) \, \mathrm{d}\alpha}{\int_0^\epsilon \frac{\alpha^t}{\alpha^{(i)}}\, \pi(\alpha) \, \mathrm{d}\alpha} \right\}\,.
\]
\end{proof}

\begin{proof}[Proof of Proposition \ref{prop:Prior}]
We first prove the upper bound. We have%On the one side
\begin{align*}
C(n,t,t+s) &\leq \frac{\int_0^\infty \frac{\alpha^{t+s}}{\alpha^{(n)}}\pi(\alpha)\, \mathrm{d}\alpha}{\int_0^\epsilon \frac{\alpha^t}{\alpha^{(n)}}\pi(\alpha) \, \mathrm{d}\alpha}
= \frac{\int_0^\epsilon \frac{\alpha^{t+s}}{\alpha^{(n)}}\pi(\alpha)\, \mathrm{d}\alpha}{\int_0^\epsilon \frac{\alpha^t}{\alpha^{(n)}}\pi(\alpha) \, \mathrm{d}\alpha}+ \frac{\int_0^\epsilon \frac{\alpha^{t+s}}{\alpha^{(n)}}\pi(\alpha)\, \mathrm{d}\alpha}{\int_0^\epsilon \frac{\alpha^t}{\alpha^{(n)}}\pi(\alpha) \, \mathrm{d}\alpha} \frac{\int_\epsilon^\infty \frac{\alpha^{t+s}}{\alpha^{(n)}}\pi(\alpha)\, \mathrm{d}\alpha}{\int_0^\epsilon \frac{\alpha^{t+s}}{\alpha^{(n)}}\pi(\alpha) \, \mathrm{d}\alpha}.
\end{align*}
Moreover, %from assumption $A2$ 
it holds 
\[
\frac{\int_\epsilon^\infty \frac{\alpha^{t+s}}{\alpha^{(n)}}\pi(\alpha)\, \mathrm{d}\alpha}{\int_0^\epsilon \frac{\alpha^{t+s}}{\alpha^{(n)}}\pi(\alpha) \, \mathrm{d}\alpha} \leq \frac{\int_\epsilon^\infty \alpha^{t+s-1}\pi(\alpha)\, \mathrm{d}\alpha}{\int_0^\epsilon \alpha^{t+s-1}\pi(\alpha) \, \mathrm{d}\alpha} \leq \delta\frac{\int_\epsilon^\infty \alpha^{t+s-1}\pi(\alpha)\, \mathrm{d}\alpha}{\int_0^\epsilon \alpha^{t+s+\beta-1}\, \mathrm{d}\alpha} \leq \delta \, E(\alpha^{t+s-1}) \frac{t+s+\beta}{\epsilon^{t+s+\beta}},
\]
where the first inequality follows since $\alpha^{(n)} \geq \epsilon^{(n)}$ for $\alpha \in (\epsilon, \infty)$ and $\alpha^{(n)} \leq \epsilon^{(n)}$ for $\alpha \in (0,\epsilon)$, while the second one follows from assumption $A2$. 
Moreover, $E$ stands for the expected value with respect to $\pi$.
Thus from Lemma \ref{main} it holds 
\[
C(n,t,t+s) \leq \frac{\delta^2 \left\{1+E(\alpha^{t+s-1})\, \frac{t+s+\beta}{\epsilon^{t+s+\beta}}  \right\}\gamma\{t+s+\beta,\epsilon\log(n)\}}{\gamma\{t+\beta,\epsilon\log(n)\}}\{\log(n)/(1+\epsilon)\}^{-s}.
\]
Then choose $G = \frac{4\delta^2}{\epsilon^{t+\beta}\gamma(t+\beta, \epsilon \log 2)}$ to obtain the upper bound. 
For the lower bound, apply Lemma \ref{main} and Lemma \ref{lower_tech} to get
\[
C(n,t,t+s) \geq \frac{1}{M+1}\frac{\int_0^\epsilon \frac{\alpha^{t+s}}{\alpha^{(n)}}\pi(\alpha)\, \mathrm{d}\alpha}{\int_0^\epsilon \frac{\alpha^t}{\alpha^{(n)}}\pi(\alpha) \, \mathrm{d}\alpha}  \geq \frac{1}{M+1}\frac{\gamma[t+s+\beta,\epsilon\{\log(n)+1\}]}{\delta^2\gamma[t+\beta,\epsilon\{\log(n)+1\}]}\{\log(n)+1\}^{-s}.
\]
Then choose $F = \frac{1}{(M+1)\delta^2 \gamma(t+\beta)}$.
\end{proof}
The following corollary of Proposition \ref{prop:Prior} will be useful.
\begin{Cor}\label{cor_C}
Suppose $\pi$ satisfies assumptions $A1$ and $A2$. Then $G >0$ as in Proposition \ref{prop:Prior} is such that for every $0 < s < n$ and $n \geq 4$ it holds
\[
\begin{aligned}
C(n,t,t+s) \leq \frac{G\Gamma(t+\beta+1)2^ss}{\epsilon}E(\alpha^{t+s-1})\log\{n/(1+\epsilon)\}^{-1}.
\end{aligned}
\]
\end{Cor}
\begin{proof}
By Proposition \ref{prop:Prior} we have
\[
C(n,t,t+s) \leq \frac{Gs}{\epsilon^s}E(\alpha^{t+s-1})\frac{\gamma\{t+s+\beta,\epsilon\log(n)\}}{\log\{n/(1+\epsilon)\}^{s}}.
\]
Note that
\[
\gamma\{t+s+\beta,\epsilon\log(n)\} = \int_0^{\epsilon \log (n)} x^{t+s+\beta-1}e^{-x}  \, \d x \leq \epsilon^{s-1}\{\log (n) \}^{s-1}\Gamma(t+\beta+1),
\]
that implies
\[
\frac{\gamma\{t+s+\beta,\epsilon\log(n)\}}{\epsilon^s \log^s\{n/(1+\epsilon)\}} \leq \frac{\Gamma(t+\beta+1)}{\epsilon}\left[\frac{\log(n)}{\log\{n/(1+\epsilon)\}} \right]^{s-1}\log\{n/(1+\epsilon)\}^{-1}.
\]
Moreover, since $\epsilon < 1$, we have $\log\{n/(1+\epsilon)\} \geq \frac{1}{2}\log(n)$ for every $n \geq 4$.
Combining the inequalities above we obtain the desired result.
\end{proof}

\section{Proof of Lemma \ref{suff_expectation}}
\begin{proof}
We need to study $R(n,1,s)$ as in \eqref{newratios}. Taking the expectation with respect to the data generating distribution we have
\begin{align*}
E\{R(n,1,s)\} &= \sum_{A \in \tau_s(n)} \frac{\prod_{j=1}^s (a_j-1)! }{(n-1)! } E\left\{\frac{\prod_{j=1}^s m(X_{A_j})}{m(X_{1:n})}\right\}\\
&=	\sum_{\bm{a} \in \mathcal{F}_s(n)} \binom{n}{a_1 \cdots a_j} \frac{\prod_{j=1}^s (a_j-1)! }{s! (n-1)! } E\left\{\frac{\prod_{j=1}^s m(X_{A_j^{\bm{a}}})}{m(X_{1:n})}\right\}\\
&=\sum_{\bm{a} \in \mathcal{F}_s(n)} \frac{n}{s! \prod_{j=1}^{s} a_j}  E\left\{\frac{\prod_{j=1}^s m(X_{A_j^{\bm{a}}})}{m(X_{1:n})}\right\}.
\end{align*}
\end{proof}

\section{Proof of Lemma \ref{priors}}
\begin{proof}
Assumptions $A1$ and $A2$ are immediately satisfied in all three cases discussed in the statement of the lemma. We thus focus on proving that $A3$ is satisfied, considering each of the three cases separately.
Suppose first that the support of the density $\pi$ is contained in $[0,c]$ with $c > 0$. Then
\[
\int_0^{\infty} \alpha^s \pi(\alpha) \, \d\alpha \leq c^s.
\]
Thus in this case assumption $A3$ is satisfied for every $\rho>0$ because $c^s < D\rho^{-s}\Gamma(s+1)$ with $D = \underset{s \in \mathbb{N}}{\max} \frac{(c\rho)^s}{\Gamma(s+1)}$ for every $\rho > 0$.
Suppose now the prior is given by a Generalized Gamma distribution, so that
\[
\int_0^{\infty} \alpha^s \pi(\alpha) \, \d\alpha = \frac{p}{a^d\Gamma\left(\frac{d}{p} \right)} \int_0^{\infty} \alpha^{d+s-1}e^{-\left(\frac{\alpha}{a} \right)^p} \, \d\alpha\,.
\]
The condition $p>1$ implies that, for every fixed $\rho > 0$ and $a > 0$, there exists $k > 0$ such that $\rho \alpha \leq \left(\frac{\alpha}{a}\right)^p$ for every $\alpha \geq k$. Thus
\[
\begin{aligned}
\int_0^{\infty} \alpha^{d+s-1}e^{-\left(\frac{\alpha}{a} \right)^p} \, \d\alpha &\leq \int_0^k \alpha^{s+d-1}e^{-\left(\frac{\alpha}{a}\right)^p} \, \d \alpha + \int_k^\infty \alpha^{s+d-1}e^{-\rho \alpha} \, \d \alpha \\
& \leq  k^{s+d-1}e^{-\left(\frac{k}{a}\right)^p} + \rho^{-d-s}\Gamma(s+d).
\end{aligned}
\]
Also,
\[
\begin{aligned}
\int_0^{\infty} \alpha^s \pi(\alpha) \, \d\alpha \leq& \frac{p}{a^d\Gamma\left(\frac{d}{p} \right)}\Gamma(s+d) \left\{  \frac{k^{s+d-1}e^{-\left(\frac{k}{a}\right)^p}}{\Gamma(s+d)} + \rho^{-d-s}\right\} \leq\\
& \leq D\rho^{-s}\Gamma(s+d),
\end{aligned}
\]
with $D = \underset{s \in \mathbb{N}}{\max} \, \frac{p}{a^d\Gamma\left(\frac{d}{p} \right)} \left\{  \frac{k^{s+d-1}e^{-\left(\frac{k}{a}\right)^p\rho^s}}{\Gamma(s+d)} + \rho^{-d}\right\}$, so that also in this case assumption $A3$ is satisfied for every $\rho > 0$. 
Finally, in the case of Gamma distribution we get
\[
\int_0^{\infty} \alpha^s \pi(\alpha) \, \d\alpha = \frac{\Gamma(\nu+s)}{\Gamma(\nu)}\rho^{-s}
\]
and assumption $A3$ holds.
\end{proof}

\section{Proof of Theorem \ref{th:ConsBounded}}

Through a linear rescaling, we may assume $[a, b] = [-c, c]$ without loss of generality.  We rewrite the assumptions on $g$ and $Q_0$ as
\begin{enumerate}
\item[$T1.$]  $\exists \, m, M$ such that $0 < m \leq g(x) \leq M < \infty$ for every $x \in [-c, c]$;
\item[$T2.$]  $g$ is differentiable on $(-c, c)$ and $\exists \, R$ such that $|\frac{g'(x)}{g(x)}| \leq R < \infty$ for every $x \in (-c, c)$;
\item[$T3.$]  $\exists \, U > 0$ such that $h(y) = q_0(y)+q_0(-y) \leq U$ for every $y \in [0, 2c]$;
\item[$T4.$]  $\exists \, L > 0$ such that $q_0(\theta) \geq L$ for every $\theta$ in a neighborhood of $\theta_j^*$, for every $j$.
\end{enumerate}
Denote with $f(x) = \sum_{j = 1}^tp_jk(x \mid \theta_j^*)$ the density of the data generating $P = \sum_{j = 1}^tp_jR_j$, with $t \in \mathbb{N}$, $p_j \in (0,1)$ and $\sum_{j = 1}^tp_j = 1$. Since $\theta^* =(\theta_1^*, \dots, \theta_t^*)$ is completely separated and \\$X^{\infty} \sim P^{(\infty)}$, each point $x$ has non-null density for at most one component of the mixture, i.e.
\[
x \in [\theta_i^*+a, \theta_i^*+b] \quad \Rightarrow \quad f(x) = p_ik(x \mid \theta_i^*) = p_ig(x-\theta_i^*).
\]
Therefore we can define
\[
C_j = \left\{ i \in \{1, \dots, n\} \, : \, x_i \in [\theta_j^*+a, \theta_j^*+b] \right\}, \quad n_j = |C_j|.
\]
Notice that $C_i \cap C_j = \emptyset$ for every $i \neq j$ and $\{1, \dots, n \} = \bigcup_{j = 1}^tC_j$, so that $\sum_{j = 1}^t n_j = n$. Moreover, defining
\[
C^{(n)} = \left\{ n_j > 0 \text{ for every $j$} \right\},
\]
for every $x_{1:n} \in C^{(n)}$ it holds
\begin{equation}\label{less_and_equal}
\begin{aligned}
&\sum_{A \in \tau_s(n)}\prod_{j=1}^s(a_j-1)!\prod_{j=1}^s m(x_{A_j}) = 0 \quad \text{for every } s < t,\\
&\sum_{B \in \tau_t(n)}\prod_{j=1}^t(b_j-1)!\prod_{j=1}^t m(x_{B_j}) = \prod_{j = 1}^t(n_j-1)!\prod_{j = 1}^tm(x_{C_j}).
\end{aligned}
\end{equation}
Since $p_j > 0$ for every $j = 1, \dots, s$, we have $P^{(n)}(C^{(n)}) \to 1$ as $n \to \infty$. We need a technical lemma.
\begin{Lem}\label{lemma:ConvProbBound}
Let $\Omega_n$ be a sequence of sets depending on $X_{1:n}$, and let $Z_n$ be random variables on the same probability space such that $P^{(\infty)}(\Omega_n) \to 1$ and
\begin{equation*}
Z_n \mathbbm{1}_{\Omega_n} \to 0
\end{equation*}
in $P^{(\infty)}$-probability as $n \to \infty$. Then $Z_n \to 0$ in $P^{(\infty)}$-probability as $n \to \infty$.
\end{Lem}
\begin{proof}
By assumption $P^{(\infty)}\left(\mathbbm{1}_{\Omega_n} Z_n > \epsilon \right) \to 0$ as $n \to \infty$. Thus, we have
\[
P^{(\infty)}\left(Z_n > \epsilon \right) \leq P^{(\infty)}\left\{(Z_n > \epsilon) \cap \Omega_n \right\}+ P^{(\infty)}\left(\Omega_n^c\right) \to 0
\]
as $n \to \infty$.
\end{proof}
Thus by Lemma \ref{lemma:ConvProbBound} it suffices to study
\begin{equation}\label{ratios_with_all_groups}
\frac{\text{pr}(K_n = s \mid X_{1:n})}{\text{pr}(K_n = t \mid X_{1:n})}\mathbbm{1}_{C^{(n)}}  = \frac{\int \dfrac{\alpha^s}{\alpha^{(n)}} \pi(\alpha) \, \mathrm{d}\alpha}{\int \dfrac{\alpha^t}{\alpha^{(n)}} \pi(\alpha) \, \mathrm{d}\alpha}\frac{\sum_{A \in \tau_s(n)}\prod_{j=1}^s(a_j-1)!\prod_{j=1}^s m(X_{A_j})}{\sum_{B \in \tau_t(n)}\prod_{j=1}^t(b_j-1)!\prod_{j=1}^t m(X_{B_j})}\mathbbm{1}_{C^{(n)}}.
\end{equation}
By \eqref{less_and_equal}, we have
\[
\frac{\text{pr}(K_n = s \mid X_{1:n})}{\text{pr}(K_n = t \mid X_{1:n})}\mathbbm{1}_{C^{(n)}}  = 0
\]
for every $s < t$. Let us now consider the case $s > t$. Again by complete separability, $A \in \tau_s(n)$ yields positive marginal density only if $A$ is a refinement of the partition $\left\{ C_1, \dots, C_t\right\}$, i.e. if
\[
A \in \tilde{\tau}_s(n) = \left\{A \in \tau_s(n) \, : \, \forall \, i = 1, \dots, s \text{ there exists } j \in \{1, \dots, t\} \text{ such that } A_i \subset C_j  \right\}.
\]
Therefore, if $A \in \tilde{\tau}_s(n)$, we write the $j$-the element as $A_j = (A^j_1, \dots, A^j_{s_j})$ with $a^j_k = |A^j_k|$, so that
\[
\sum_{A \in \tilde{\tau}_s(n)}\prod_{j=1}^s(a_j-1)!\prod_{j=1}^s m(X_{A_j}) = \sum_{\textbf{s}\in \text{S}}\prod_{j = 1}^t\sum_{A_j \in \tau_{s_j}(n_j)}\prod_{k=1}^{s_j}(a^j_k-1)!\prod_{k=1}^{s_j} m(X_{A^j_k}),
\]
where $\textbf{S} = \left\{(s_1, \dots, s_t) \, : \, 1 \leq s_j \leq n_j, \, \forall j, \text{ and } \sum_{j = 1}^ts_j = s \right\}$. By the above and \eqref{less_and_equal} we can rewrite \eqref{ratios_with_all_groups} as
\begin{equation}\label{ratios_for_inconsistency}
\begin{aligned}
\frac{\text{pr}(K_n = s \mid X_{1:n})}{\text{pr}(K_n = t \mid X_{1:n})}\mathbbm{1}_{C^{(n)}} & = C(n, t, s)\frac{\sum_{A \in \tilde{\tau}_s(n)}\prod_{j=1}^s(a_j-1)!\prod_{j=1}^s m(X_{A_j})}{\prod_{j = 1}^t(n_j-1)!\prod_{j = 1}^tm(X_{C_j})}\mathbbm{1}_{C^{(n)}} \\
&= C(n, t, s)\sum_{\textbf{s}}\prod_{j = 1}^t\sum_{A_j \in \tau_{s_j}(n_j)}\frac{\prod_{k=1}^{s_j}(a^j_k-1)!}{(n_j-1)!}\frac{\prod_{k=1}^{s_j} m(X_{A^j_k})}{m(A_{C_j})}\mathbbm{1}_{C^{(n)}},
\end{aligned}
\end{equation}
where
\[
m(X_{C_j}) = \int_{\mathbb{R}}\prod_{i \in C_j}k(X_i \mid \theta_j) \, Q_0(\d \theta_j) = \int_{\mathbb{R}}\prod_{i \in C_j}g(X_i - \theta_j) \, Q_0(\d \theta_j)
\]
and
\[
m(X_{A^j_h}) = \int_{\mathbb{R}}\prod_{i \in A^j_h}k(X_i \mid \theta_h) \, Q_0(\d \theta_h) = \int_{\mathbb{R}}\prod_{i \in A^j_h}g(X_i - \theta_h) \, Q_0(\d \theta_h),
\]
with $h = 1, \dots, s_j$. We divide and multiply by
\[
\prod_{i = 1}^nf(X_i) = \prod_{j = 1}^t\prod_{i \in C_j}p_jk(X_i \mid \theta_j^*) = \prod_{j = 1}^t\prod_{h = 1}^{s_j}\prod_{i \in A_h^j}p_jk(X_i \mid \theta_j^*),
\]
so that the sum on the right hand side of \eqref{ratios_for_inconsistency} becomes
\begin{equation}\label{object}
\sum_{\textbf{s}}\prod_{j = 1}^t\sum_{A_j \in \tau_{s_j}(n_j)}\frac{\prod_{k=1}^{s_j}(a^j_k-1)!}{(n_j-1)!}\frac{\prod_{k = 1}^{s_j}\int_{\mathbb{R}}\prod_{i \in A^j_k}\frac{g(X_i - \theta_k)}{p_jg(X_i-\theta_j^*)} \, Q_0(\d \theta_k)}{\int_{\mathbb{R}}\prod_{i \in C_j}\frac{g(X_i - \theta_j)}{p_jg(X_i-\theta_j^*)} \, Q_0(\d \theta_j)}\mathbbm{1}_{C^{(n)}}, \quad \text{for } s > t.
\end{equation}
We start with the denominator. The next lemma specifies the behaviour of the maximum for each group, where $X^j_{(r)}$ denotes the $r$-th order statistic of $X_{C_j}$.
\begin{Lem}\label{lemma:Max}
For every $j = 1, \dots , t$ it holds
\[
Y^j_{n_j} := \min \left[ 1, n_j(\log(n))^{\frac{1}{2t}}\{c+\theta_j^*-X^j_{(n_j)}\} \right] \to 1
\]
in $P^{(\infty)}$-probability as $n \to \infty$.
\end{Lem}
\begin{proof}
First, notice that $n_j \to \infty$ $P^{(\infty)}$-almost surely as $n \to \infty$. By definition $Y^j_{n_j} \leq 1$, so we have to prove that $\forall \epsilon > 0$
\[
P^{(\infty)}\left(1- Y^j_{n_j} > \epsilon \right) \to 0
\]
as $n_j \to \infty$, where pr is evaluated with respect to $P^{(\infty)}$. Without loss of generality assume $\theta_j^* = 0$. Thus, by definition we have
\[
\begin{aligned}
P^{(\infty)}(1 - Y^j_{n_j} > \epsilon) &=P^{(\infty)}\left[n_j(\log(n))^{\frac{1}{2t}}\{c-X^j_{(n)} \} \leq 1-\epsilon \right] = P^{(\infty)}\left\{ X^j_{(n)} \geq c-\frac{1-\epsilon}{n_j(\log(n))^{\frac{1}{2t}}}\right\} \\
&= 1-\left\{1-\int_{c-\frac{1-\epsilon}{n_j(\log(n))^{\frac{1}{2t}}}}^c g(x) \, \mathrm{d}x \right\}^{n}\,.
\end{aligned}
\]
Thus, by $T1$ we have that $\int_{c-\frac{1-\epsilon}{n_j(\log(n))^{\frac{1}{2t}}}}^c g(x) \, \mathrm{d}x \leq \frac{M(1-\epsilon)}{n_j(\log(n))^{\frac{1}{2t}}}$, so that
\[
P^{(\infty)}(1 - Y^j_{n_j} > \epsilon) \leq 1-\left\{ 1- \frac{M(1-\epsilon)}{n_j(\log(n))^{\frac{1}{2t}}}\right\}^{n}= 1-e^{-\frac{M(1-\epsilon)}{(\log(n))^{\frac{1}{2t}}}+n_j \, \text{o}\left(\frac{1}{n_j(\log(n))^{\frac{1}{2t}}}\right)} \to 0,
\]
as $n \to \infty$, by the Taylor expansion of the logarithmic function.
\end{proof}
\begin{Lem}\label{denom1}
For every $j = 1, \dots, t$ it holds
\[
\prod_{i \in C_j} \frac{g(x_i - \theta_j)}{g(x_i)} \geq e^{-R} \mathbbm{1}_{[0,\frac{1}{n_j}]} (|\theta_j-\theta_j^*|) \mathbbm{1}_{[x^j_{(n_j)}-c,x^j_{(1)}+c]}(\theta_j-\theta_j^*).
\] 
with $R$ defined in $T2$ and $x^j_{(r)}$ denotes the $r$-th order statistic of $x_{C_j}$.
\end{Lem}
\begin{proof}
Without loss of generality assume $\theta_j^* = 0$. Define $p(x) := \log g(x)$, with $x \in [-c, c]$, so that $p'(x) = \frac{g'(x)}{g(x)}$. By $T2$ and the Fundamental Theorem of Integral Calculus
\[
|p(y)-p(x)| = \left\lvert \int_x^yp'(t) \, \d t \right\rvert \leq \int_x^y \left|\frac{g'(t)}{g(t)}\right| \d t \leq R|y-x|, \quad -c < x \leq y < c.
\]
Thus, we have 
\[
\frac{g(x - \theta_j)}{g(x)} = e^{p(x -\theta_j)-p(x)} = e^{-\{p(x)-p(x - \theta_j)\}} \geq e^{-R|\theta_j|}, \quad x \in [-c,c].
\]
Finally, we get
\[
\begin{aligned}
\prod_{i \in C_j} \frac{g(x_i - \theta_j)}{g(x_i)} &\geq e^{-Rn_j|\theta_j|}\mathbbm{1}_{[x^j_{(n_j)}-c,x^j_{(1)}+c]}(\theta_j) \geq e^{-Rn|\theta_j|}\mathbbm{1}_{[0,\frac{1}{n_j}]} (|\theta_j|)\mathbbm{1}_{[x^j_{(n_j)}-c,x^j_{(1)}+c]}(\theta_j)\\
& \geq e^{-R} \mathbbm{1}_{[0,\frac{1}{n_j}]} (|\theta_j|) \mathbbm{1}_{[x^j_{(n_j)}-c,x^j_{(1)}+c]}(\theta_j).
\end{aligned}
\]
\end{proof}
\begin{Lem}\label{denom}
For every $j = 1, \dots, t$ there exists $K>0$ and $N_j \in \mathbb{N}$ such that for all $n_j\geq N_j$  it holds
\[
\int_\R \prod_{i \in C_j}\frac{g(X_i - \theta_j)}{g(X_i-\theta_j^*)}q_0(\theta_j) \, \mathrm{d}\theta_j %&\geq Le^{-2c}\min \left\{\frac{1}{n}, c-X_{(n)} \right\} = \frac{Le^{-2c}\min \left\{1, n(c-X_{(n)}) \right\}}{n} \geq \\
\geq %\frac{c \min \left[1, n\sqrt{\log(n)}\{c-X_{(n)}\} \right]}{n\sqrt{\log(n)}} = 
\frac{K^{\frac{1}{t}} Y^j_{n_j} }{n_j(\log(n))^{\frac{1}{2t}}},
\]
with $Y^j_{n_j}$ defined in Lemma \ref{lemma:Max}.
\end{Lem}
\begin{proof}
Without loss of generality assume $\theta_j^* = 0$. Notice that, by $T4$, there exists $N_j \in \mathbb{N}$ such that $q_0(\theta) \geq L$ for every $\theta \in \left[-\frac{1}{N_j},0\right]$. Thus, applying Lemma \ref{denom1} and considering $n_j \geq N_j$, we get
\[
\begin{aligned}
\int_\R \prod_{i \in C_j}\frac{g(X_i - \theta_j)}{g(X_i)}q_0(\theta_j) \, \mathrm{d}\theta_j&\geq e^{-R}\int_\R \mathbbm{1}_{[0,\frac{1}{n_j}]} (|\theta_j|) \mathbbm{1}_{[X^j_{(n_j)}-c,x^j_{(1)}+c]}(\theta_j) \, q_0(\theta_j)\, \mathrm{d}\theta_j\\
%& = e^{-2c}\int_{-\frac{1}{n}}^{\frac{1}{n}}\mathbbm{1}_{[x_{(n)}-c,x_{(1)}+c]}(\theta) \, q_0(\theta) \, \mathrm{d}\theta \geq \\
& \geq e^{-R}\int_{-\frac{1}{n_j}}^0\mathbbm{1}_{\{X^j_{(n_j)} \leq \theta_j+c\}} \, q_0(\theta_j) \, \mathrm{d}\theta_j \geq Le^{-R}\min \left\{\frac{1}{n_j}, c-X^j_{(n_j)} \right\},
\end{aligned}
\]
with $L$ defined in $T4$. Thus, multiplying both the numerator and the denominator by $n_j(\log(n))^{\frac{1}{2t}}$, with $n \geq N$, we have
\[
\begin{aligned}
\int_\R \prod_{i \in C_j}\frac{g(X_i - \theta_j)}{g(X_i)}q_0(\theta_j) \, \mathrm{d}\theta_j &\geq 2 Le^{-R}\min \left\{\frac{1}{n_j}, c-X^j_{(n_j)} \right\} \\%= \frac{c \min \left\{1, n(c-X_{(n)}) \right\}}{n} \geq \\
& \geq \frac{K^{\frac{1}{t}} \min \left[1, n_j(\log(n))^{\frac{1}{2t}}\{c-X_{(n)}\} \right]}{n_j(\log(n))^{\frac{1}{2t}}} =  \frac{K^{\frac{1}{2t}} Y_n }{n_j(\log(n))^{\frac{1}{2t}}},
\end{aligned}
\]
with $K = (2Le^{-R})^t$.
\end{proof}
Define the event
\begin{equation}\label{Omega_n}
\Omega_n = \left\{\text{for every $j = 1, \dots, t$ it holds: } n_j \geq N_j, Y^j_{n_j} \in [1/2, 1] \right\},
\end{equation}
such that $P^{(n)}(\Omega_n) \to 1$ thanks to Lemma \ref{lemma:Max} and Lemma \ref{denom}. Thus, an upper bound of \eqref{object} with $\Omega_n$ in place of $C^{(n)}$ is given by
\begin{equation}\label{upper_bound_ratios}
T^{(n)} := \frac{2^t\sqrt{\log(n)}}{K}\sum_{\textbf{s}}\prod_{j = 1}^t\sum_{A_j \in \tau_{s_j}(n_j)}n_j\frac{\prod_{k=1}^{s_j}(a^j_k-1)!}{(n_j-1)!}\prod_{h = 1}^{s_j}\int_{\mathbb{R}}\prod_{i \in A^j_h}\frac{g(X_i - \theta_h)}{g(X_i-\theta_j^*)} \, Q_0(\d \theta_h)\mathbbm{1}_{\Omega_n},
\end{equation}
for $s > t$. Now we apply the expected value with respect to the values of each group, as shown in the next lemma.

\begin{Lem}\label{numerator}
Under $X_{1:n} \sim P^{(n)}$, for every $j = 1, \dots, t$, $s_j \geq 1$ and $(\theta_1, \dots, \theta_{s_j}) \in \mathbb{R}^{s_j}$, we have
\[
\begin{aligned}
E \left\{\prod_{h = 1}^{s_j}\int_{\R^{s_j}} \prod_{i \in A^j_h} \frac{g(X_i - \theta_h)}{g(X_i-\theta_j^*)} q_0(\theta_h)\, \mathrm{d}\theta_h \right\} 
&\leq \left(\frac{U}{m}\right)^{s_j}\prod_{h= 1}^{s_j} \frac{1}{a^j_h+1},
\end{aligned}
\]
with $m$ and $U$ defined in $T1$ and $T3$. 
\end{Lem}
\begin{proof}
Without loss of generality assume $\theta_j^* = 0$. Taking the expectation under $P^{(n)}$ we have
\begin{equation}\label{expectation_num}
\begin{aligned}
E \left\{\int_{\R^{s_j}} \prod_{h = 1}^{s_j} \prod_{i \in A^j_h} \frac{g(X_i - \theta_h)}{g(X_i)} q_0(\theta_h)\, \mathrm{d}\theta_h \right\} %&= \int_{[-c, c]^n}\int_{\Theta^s} \prod_{j = 1}^s \prod_{i \in A_j} k(x_i - \theta_j) q_0(\theta_j) \, \mathrm{d}\theta_j\, \mathrm{d}x_i\\
& = \int_{\R^{s_j}} \int_{[-c,c]^{n_j}} \prod_{h = 1}^{s_j} \prod_{i \in A^j_h} g(x_i - \theta_h) q_0(\theta_h) \, \mathrm{d}x_i \, \mathrm{d}\theta_h,
\end{aligned}
\end{equation}
By the change of variables $z = x-\theta_h$, we have
\[
\int_{-c}^{c}  g(x - \theta_h) \mathbbm{1}_{[\theta_h-c, \theta_h+c]}(x)\, \mathrm{d}x = \int_{-c-\theta_h}^{c-\theta_h} g(z)\mathbbm{1}_{[-c, c]}(z)\, \mathrm{d}z.
\]
If $\theta_h > 0$, then
\[
\begin{aligned}
\int_{-c-\theta_h}^{c-\theta_h} g(z)\mathbbm{1}_{[-c, c]}(z)\, \mathrm{d}z &= \mathbbm{1}_{[0, 2c]}(\theta_h)\int_{-c}^{c-\theta_h}g(z) \, \mathrm{d}z\\
& = \mathbbm{1}_{[0, 2c]}(\theta_h)\left(1-\int_{c-\theta_h}^c g(z) \, \mathrm{d}z\right) \leq \mathbbm{1}_{[0, 2c]}(|\theta_h|)\left(1-m|\theta_h|\right).
\end{aligned}
\]
Similarly, if $\theta_h < 0$ we get
\[
\begin{aligned}
\int_{-c-\theta_h}^{c-\theta_h} g(z)\mathbbm{1}_{[-c, c]}(z)\, \mathrm{d}z &= \mathbbm{1}_{[-2c, 0]}(\theta_h)\int_{-c-\theta_h}^{c}g(z) \, \mathrm{d}z\\
&= \mathbbm{1}_{[-2c, 0]}(\theta_h)\left( 1-\int_{-c}^{-c-\theta_h} g(z) \, \mathrm{d}z\right) \leq \mathbbm{1}_{[0,2c]}(|\theta_h|) \left(1-m|\theta_h|\right).
\end{aligned}
\]
Thus
\[
\int_{-c}^{c}  g(x - \theta_h) \mathbbm{1}_{[\theta_h-c, \theta_h+c]}(x)\, \mathrm{d}x  \leq \mathbbm{1}_{[0,2c]}(|\theta_h|) \left(1-m|\theta_h|\right), \quad h = 1, \dots, s_j,
\]
which implies
\[
\prod_{h = 1}^{s_j} \prod_{i \in A^j_h} \int_{-c}^{c}  g(x - \theta_h) \mathbbm{1}_{[\theta_h-c, \theta_h+c]}(x)\, \mathrm{d}x  \leq \prod_{h = 1}^{s_j}\mathbbm{1}_{[0,2c]}(|\theta_h|) \left(1-m|\theta_h|\right).
\]
Considering $h$ defined as in $T3$, we have
\[
\int_{\R}\mathbbm{1}_{[0,2c]}(|\theta_h|) \left(1-m|\theta_h|\right)q_0(\theta_h) \, \d \theta_h = \int_0^{2c} \left(1-m|\theta_h|\right)h(\theta_h) \, \d \theta_h, \quad h =1, \dots, s_j.
\]
Combining the above with \eqref{expectation_num} we get
\begin{equation}\label{expectation_num2}
\begin{aligned}
E \left\{\int_{\R^{s_j}} \prod_{h = 1}^{s_j} \prod_{i \in A^j_h} \frac{g(X_i - \theta_h)}{g(X_i)} q_0(\theta_h)\, \mathrm{d}\theta_h \right\} %&= \int_{[-c, c]^n}\int_{\Theta^s} \prod_{j = 1}^s \prod_{i \in A_j} k(x_i - \theta_j) q_0(\theta_j) \, \mathrm{d}\theta_j\, \mathrm{d}x_i\\
& = \int_{\R^{s_j}} \int_{[-c,c]^{n_j}} \prod_{h = 1}^{s_j} \prod_{i \in A^j_h} g(x_i - \theta_h) q_0(\theta_h) \, \mathrm{d}x_i \, \mathrm{d}\theta_h\\
&\leq \prod_{h=1}^{s_j}\int_0^{2c} \left(1-m|\theta_h|\right)h(\theta_h) \, \d \theta_h.
\end{aligned}
\end{equation}
With $U$  defined as in $T3$, we have
\[
\int_{0}^{2c} (1-my)^{a_h^j} h(y) \, \mathrm{d}y \leq U \int_{0}^{2c} (1-m y)^{a_h^j} \, \mathrm{d}y.
\]
Now consider the change of variables $u = 1-my$ and compute
\[
\begin{aligned}
\int_{0}^{2c} (1-my)^{a_h^j} \, \mathrm{d}y &= \frac{1}{m} \int_{1-2 m c}^1u^{a_h^j} \, \mathrm{d}u = \frac{1-(1-2mc)^{a_h^j+1}}{m(a_h^j+1)} \leq \frac{1}{m(a_h^j+1)}.
\end{aligned}
\]
Finally, through \eqref{expectation_num2}, we have
\[
\begin{aligned}
E \left\{\int_{\R^{s_j}} \prod_{h = 1}^{s_j} \prod_{i \in A^j_h} \frac{g(X_i - \theta_h)}{g(X_i)} q_0(\theta_h)\, \mathrm{d}\theta_h \right\}  %&= \int_{[-c, c]^n}\int_{\Theta^s} \prod_{j = 1}^s \prod_{i \in A_j} k(x_i - \theta_j) q_0(\theta_j) \, \mathrm{d}\theta_j\, \mathrm{d}x_i\\
& \leq \prod_{h=1}^{s_j}\int_0^{2c} \left(1-m|\theta_h|\right)h(\theta_h) \, \d \theta_h\\
&\leq \left(\frac{U}{m}\right)^{s_j}\prod_{h = 1}^{s_j} \frac{1}{a_h^j+1},
\end{aligned}
\]
as desired.
\end{proof}


\subsection{Proof of Theorem \ref{th:ConsBounded}}
We have the next two technical lemmas.
\begin{Lem}\label{bound_sum}
Let $p^* = \min_{j \in \{1, \dots, t\}} p_j \in (0,1)$. It holds
\[
\sum_{\textbf{s} \in \textbf{S}}\frac{s!}{\prod_{j = 1}^ts_j!} = \sum_{\textbf{s}}\binom{s}{s_1, \dots, s_t} \leq (p^*)^{-s},
\]
where $\textbf{S} = \left\{(s_1, \dots, s_t) \, : \, s_j \leq n_j \text{ and } \sum_{j = 1}^ts_j = s \right\}$.
\end{Lem}
\begin{proof}
The result follows immediately from
\[
\begin{aligned}
\sum_{\textbf{s} \in \textbf{S}}\binom{s}{s_1, \dots, s_t} &\leq (p^*)^{-s}\sum_{\textbf{s} \in \textbf{S}}\binom{s}{s_1, \dots, s_t}\prod_{j = 1}^tp_j^{s_j}\\
& \leq (p^*)^{-s}\sum_{\textbf{s} \in R_t}\binom{s}{s_1, \dots, s_t}\prod_{j = 1}^tp_j^{s_j}= (p^*)^{-s},
\end{aligned}
\]
where $R_t = \left\{(s_1, \dots, s_t) \, : \, \sum_{j = 1}^ts_j = s \right\}$, since the sum on the right-hand side is the sum of the probabilities over all the possible values of a multinomial distribution with parameters $(s, p_1, \dots, p_t)$.
\end{proof}
\begin{Lem}\label{induction}
For every $p>1$ and for every integers $s\geq 2$ and $n\ge s$ it holds
\begin{align*}
\sum_{\bm{a} \in \mathcal{F}_s(n)} \left(\frac{n}{\prod_{j=1}^s a_j}\right)^p < C_p^{s-1},
\end{align*}
where $\mathcal{F}_s(n) = \left\{ \textbf{a} \in \{1,\ldots,n\}^s: \sum_{j = 1}^s a_j =n\right\}$
and  $C_p = 2^{p}\zeta(p)$, with $\zeta(p) = \sum_{a=1}^{\infty} \frac{1}{a^p}<\infty$.
\end{Lem}
\begin{proof}
We prove the result by induction. Consider the base case $s=2$. 
By the strict convexity of $x\mapsto x^p$ for $p>1$ we have
\begin{align*}
\sum_{\bm{a} \in \mathcal{F}_2(n)} \left(\frac{n}{a_1 a_2}\right)^p = \sum_{a=1}^{n-1} \left\{\frac{n}{a(n-a)}\right\}^p
= 2^p \sum_{a=1}^{n-1} \left(\frac{1}{2}\frac{1}{a} +\frac{1}{2} \frac{1}{n-a}\right)^p < 2^{p} \sum_{a=1}^{n-1} \frac{1}{a^p} < C_p,
\end{align*}
for every $n\geq 2$. 
For the induction step, assume that for some $s\geq 3$ we have
$$
\sum_{\bm{a}\in \mathcal{F}_{s-1}(n)} \left(\frac{n}{\prod_{j=1}^{s-1} a_j}\right)^2 < C_p^{s-2}
$$
for all $n\geq s-1$. Then
\begin{align*}
\sum_{\bm{a} \in \mathcal{F}_s(n)} \left(\frac{n}{\prod_{j=1}^s a_j}\right)^p 
&=
\sum_{a_{s}=1}^{n-s+1} \sum_{(a_1,\dots,a_{s-1})\in \mathcal{F}_{s-1}(n-a_{s})} \left(\frac{n}{\prod_{j=1}^s a_j}\right)^p 
\\&=
\sum_{a_{s}=1}^{n-s+1} \left\{\frac{n}{(n-a_{s})a_{s}}\right\}^p \sum_{(a_1,\dots,a_{s-1})\in \mathcal{F}_{s-1}(n-a_{s})} \left(\frac{n-a_s}{\prod_{j=1}^{s-1} a_j}\right)^p\\
&\le C_p^{s-2} \sum_{a_{s}=1}^{n-s+1} \left\{\frac{n}{(n-a_{s})a_{s}}\right\}^p  < C_p^{s-1}
\end{align*}
%meaning that the desired statement holds for $s$. T
and thus the thesis follows by induction.
\end{proof}
In the following we will drop the subscript in $C_p$ when the value of $p$ is clear from the context, thus denoting $C = C_p$.
\begin{Lem}\label{lemma:consGeneral}
Consider the setting of \eqref{eq:MixDPM} with $(f, k, q_0)$ as in Theorem \ref{th:ConsBounded}. Moreover, assume $\pi(\alpha)$ satisfies assumptions $A1$, $A2$, and $A3$. Then, under $X_{1:\infty} \sim P^{(\infty)}$ we have
\[
E\left\{\mathbbm{1}_{\Omega_n} \sum_{s = 1}^{n-t} \frac{\text{pr}(K_{n} = t+s \, | \, X_{1:n})}{\text{pr}(K_{n} = t \, | \, X_{1:n})} \right\} \to 0
\]
as $n \to \infty$, with $\Omega_n$ as in \eqref{Omega_n}.
\end{Lem}
\begin{proof}
Applying Lemma \ref{numerator} we can upper bound the expected value of $T^{(n)}$ in \eqref{upper_bound_ratios} as follows
\[
\begin{aligned}
\mathbb{E}\left\{T^{(n)}\right\} \leq \frac{2^t\sqrt{\log(n)}}{K}\left(\frac{U}{m}\right)^s&\sum_{\textbf{s}}\prod_{j = 1}^t\sum_{A_j \in \tau_{s_j}(n_j)}\frac{n_j}{(n_j-1)!\prod_{k = 1}^{s_j}(a^j_k+1)}\\
& \leq \frac{2^t\sqrt{\log(n)}}{K}\left(\frac{U}{m}\right)^s\sum_{\textbf{s}}\prod_{j = 1}^t\frac{1}{s_j!}\sum_{\textbf{a}_j \in \mathcal{F}_{s_j}(n_j)}\left(\frac{n_j}{\prod_{k = 1}^{s_j}a^j_k}\right)^2,
\end{aligned}
\]
where the last inequality follows from Lemma \ref{suff_expectation}. Moreover, from Lemma \ref{induction} we have
\[
\sum_{\textbf{a}_j \in \mathcal{F}_{s_j}(n_j)}\left(\frac{n_j}{\prod_{k = 1}^{s_j}a^j_k}\right)^2 < C^{s_j},
\]
with constant $C < 7$. Thus
\begin{equation}\label{bound_T}
\mathbb{E}\left\{T^{(n)}\right\} \leq \frac{2^t\sqrt{\log(n)}}{K}\left(\frac{UC}{m}\right)^s\sum_{\textbf{s}}\prod_{j = 1}^t\frac{1}{s_j!}.
\end{equation}
Moreover, from Corollary \ref{cor_C} and $A3$ we have
\begin{equation}\label{bound_C}
\begin{aligned}
C(n,t,t+s) &\leq \frac{G\Gamma(t+\beta+1)2^ss}{\epsilon}E(\alpha^{t+s-})\log\{n/(1+\epsilon)\}^{-1} \\
& \leq \frac{DG\Gamma(t+\beta+1)2^ss}{\epsilon}\rho^{-(t+s-1)}\Gamma(\nu+t+s)\log\{n/(1+\epsilon)\}^{-1}, \quad n \geq 4\,.
\end{aligned}
\end{equation}
By \eqref{bound_T}, combined with Lemma \ref{bound_sum}, and \eqref{bound_C} we finally have
\[
\begin{aligned}
E \biggl\{\mathbbm{1}_{\Omega_n}&\sum_{s = 1}^{n-t}\frac{\text{pr}(K_n = s+t | X_{1:n})}{\text{pr}(K_n = t | X_{1:n})}\biggr\} = \sum_{s = 1}^{n-t}C(n,t,t+s)E\{\mathbbm{1}_{\Omega_n}R(n,t,t+s)\}\\
&  \leq \frac{2^t\rho^{1-t}(U/m)^tDG\Gamma(t+\beta+1)\sqrt{\log(n)}}{K\epsilon \log\{n/(1+\epsilon)\}}\underbrace{\sum_{ s = 1}^{n-1}\frac{s (2CUp^*/m)^{s}\rho^{-s}\Gamma(\nu +t+s)}{(s+1)!}}_{< \infty} \to 0,
\end{aligned}
\]
as $n \to \infty$, where finiteness follows by taking $\rho$ sufficiently large.
\end{proof}
\begin{proof}[Proof of Theorem \ref{th:ConsBounded}] First of all, assume $\pi(\cdot)$ satisfies $A1-A3$. By Lemma \ref{lemma:consGeneral} it holds
\[
\mathbbm{1}_{\Omega_n}\sum_{s = 1}^{n-t} \frac{\text{pr}(K_{n} = t+s \, | \, X_{1:n})}{\text{pr}(K_{n} = t \, | \, X_{1:n})} \to 0
\]
in $P^{(\infty)}$--probability as $n \to \infty$. The desired result then follows from Lemma \ref{lemma:ConvProbBound} with $Z_n = \sum_{s = 1}^{n-t} \frac{\text{pr}(K_{n} = t+s \, | \, X_{1:n})}{\text{pr}(K_{n} =t \, | \, X_{1:n})}$ and $\Omega_n$ as in \eqref{Omega_n}.

Assume instead $\pi(\alpha) = \delta_{\alpha^*}(\alpha)$ with $\alpha^* > 0$. By \eqref{ratios_for_inconsistency} we have
\[
\frac{p(K_n = t+1 \mid X_{1:n})}{p(K_n = t \mid X_{1:n})} \geq \alpha^*\sum_{\textbf{s} \in \textbf{S}}\prod_{j = 1}^t\sum_{A_j \in \tau_{s_j}(n_j)}\frac{\prod_{k=1}^{s_j}(a^j_k-1)!}{(n_j-1)!}\frac{\prod_{k=1}^{s_j} m(X_{A^j_k})}{m(A_{C_j})}.
\]
Notice that, with $n$ high enough, $n_1 > 1$ almost surely. Then, denoting $i \in C_1$, we consider the special case
\[
\textbf{s} = (2,1, \dots, 1), \quad A^1_1 = \{i\}, A_2^1 = A_{C_1} \backslash \{i\},
\]
and $A_j = \{A_{C_j}\}$ for every $j \geq 2$. Thus we can write
\begin{equation}\label{inequality_inconsistency}
\frac{p(K_n = t+1 \mid X_{1:n})}{p(K_n = t \mid X_{1:n})} \geq \alpha^*\sum_{i \in C_1}\frac{1}{n_1-1}\frac{m(X_i)m\left(X_{C_1 \backslash i} \right)}{m\left(X_{C_j} \right)}.
\end{equation}
By $T1$ we have
\[
\begin{aligned}
m\left(X_{C_j} \right) &= \int_{\R} \prod_{j \in C_1}g(X_j-\theta)q_0(\theta) \, \d \theta\\
& \leq M \int_{\R} \prod_{j \in C_1 \backslash i}g(X_j-\theta)q_0(\theta) \, \d \theta = M \, m\left(X_{C_1 \backslash i} \right).
\end{aligned}
\]
Moreover, by $T4$ there exists $\epsilon > 0$ such that
\[
\begin{aligned}
m(X_i) = \int_{\R} g(X_i -\theta) q_0(\theta) \d \theta &\geq m \int_{\theta_1^*-\epsilon}^{\theta_1^* + \epsilon}q_0(\theta) \d \theta
&\geq 2mL\epsilon.
\end{aligned}
\]
Therefore, \eqref{inequality_inconsistency} becomes
\[
\frac{p(K_n = t+1 \mid X_{1:n})}{p(K_n = t \mid X_{1:n})} \geq \frac{2\alpha^*mL\epsilon}{M}\sum_{i \in C_1}\frac{1}{n_1-1} = \frac{2\alpha^*mL\epsilon}{M}\frac{n_1}{n_1-1},
\]
and
\[
\lim \inf_{n \to \infty} \sum_{s \neq t}\frac{p(K_n = s \mid X_{1:n})}{p(K_n = t \mid X_{1:n})} \geq \lim \inf_{n \to \infty} \frac{p(K_n = t+1 \mid X_{1:n})}{p(K_n = t \mid X_{1:n})} \geq \frac{\alpha^*mL\epsilon}{M} > 0.
\]
Then
\[
\begin{aligned}
\lim \sup_{n \to \infty} \text{pr}(K_n = t \mid X_{1:n}) &= \lim \sup_{n \to \infty} \left\{1+\sum_{s \neq t} \frac{\text{pr}(K_n = s \mid X_{1:n} )}{\text{pr}(K_n = t \mid X_{1:n})}\right\}^{-1}\\
& = \frac{1}{1+\lim \inf_{n \to \infty} \sum_{s \neq t}\frac{p(K_n = s \mid X_{1:n})}{p(K_n = t \mid X_{1:n})}} > 0,
\end{aligned}
\]
which completes the proof.
\end{proof}

\section{Proof of Proposition \ref{undersmoothing}}
We adapt the proof of Theorem $2.1$ in \cite{Cai2020}. Denote by
\[
\Psi = \left\{k(\cdot \mid \theta): \theta \in \Theta \subseteq \mathbb{R}^p \right\}
\]
the family of kernels, dominated by $\mu$, either Lebesgue or counting measure, and with common domain $\mathbb{X} \subseteq \mathbb{R}^q$. Denote with $B_x(\epsilon)$ the closed ball of center $x \in \mathbb{X}$ and radius $\epsilon > 0$. Let $\bar{\Theta}$ be the closure of $\Theta$ and define the set
\[
\mathbb{B} := \left\{ \bar{\theta} \in \bar{\Theta}\backslash \Theta \, : \, \lim_{\theta \to \bar{\theta}} \left \{\sup_x k(x \mid \theta)\right\} = \infty \right\}.
\]
Let $\mathbb{G}_s$ be the set of mixtures of exactly $s$ elements in $\Psi$, that is
\[
f \in \mathbb{G}_s \quad \Leftrightarrow \quad f = \sum_{j = 1}^sq_jk(\cdot \mid \theta_j),
\]
with $q_j > 0$ for every $j$, $\sum_{j = 1}^sq_j = 1$ and $\theta_i \neq \theta_h$ for every $i \neq h$.  Let $\mathcal{P}(G)$ be the set of probability measures on a generic space $G$; with a slight abuse of notation we will say $f \in \mathcal{P}(G)$ when $f$ is the density of a probability measure $P \in \mathcal{P}(G)$. Therefore, given $P \in \mathbb{G}_t$, with weights $\{p_j\}_{j = 1}^t$ and parameters $\{\theta_j^* \}_{j = 1}^t$, we define the Kullback-Leibler neighborhoods of $P$ as
\begin{equation}\label{KL_neighborhood}
KL_\epsilon(P) := \bigg\{ h \in \mathcal{P}(\mathbb{X}): \int  \log\bigg\{ \frac{\sum_{j=1}^{t} p_j k(x\mid \theta_j^\ast)}{h(x)} \bigg\} P(\mathrm{d}x)  < \epsilon \bigg\},
\end{equation}
for $\epsilon > 0$. We make the following assumptions: 
\begin{enumerate}
\item[$H1.$] For every $\bar{\theta} \in \Theta \backslash \mathbb{B}$, for $\mu$-almost every $x \in \mathbb{X}$ there exists $A := A(\bar{\theta}, x) \subset \Theta \backslash \mathbb{B}$ neighborhood of $\bar{\theta}$ so that the mapping $\theta \in A \to k(x \mid \theta)$ is continuous. Moreover $\mathbb{B}$ is closed; 
\item[$H2.$] Let $\{\theta_i \}_{i = 1}^\infty \subset \Theta$. If $||\theta_i|| \to \infty$ as $i \to \infty$, then for every compact set $K \subset \mathbb{X},$ \\$\int_Kk(x\mid \theta_i) \, \mu(\d x) \to 0$, as $i \to \infty$. If $\theta_i \to \bar{\theta} \in \mathbb{B}$, then there exists $x^* \in \mathbb{X}$ such that $k(\cdot \mid \theta_i) \to \delta_{x^*}(\cdot)$ weakly as $i \to \infty$;% for every $\epsilon > 0$ \\$\int_{B_{x^*}(\epsilon)}k(x\mid \theta_i) \, \mu(\d x) \to 1$, as $i \to \infty$;
\item[$H3.$] If $f \in \mathbb{G}_t$, then there exist no $f' \in \mathbb{G}_s$, with $s < t$, such that $f(x) = f'(x)$ $\mu$-almost surely;
\item[$H4.$] For every $P \in \mathbb{G}_t$, $t \geq 1$, with $\theta_1^*, \dots, \theta_t^*$ belonging to the support of $Q_0$, we have \\$\text{pr}(h \in K_\epsilon(P))>0$ for every $\epsilon > 0$, where $h$ follows the prior distribution in \eqref{eq:MixDPM}.
\end{enumerate}
Assumption $H2$ says that, when $\theta$ diverges or converges to elements in $\mathbb{B}$, the kernel $k$ degenerates: it is satisfied for instance when the elements of $\theta$ are location or scale parameters. $H3$ instead implies that the clustering problem is not ill-posed, in the sense that different numbers of components always lead to different distribution. $H4$ finally requires that the finite mixtures of the kernel $k(\cdot \mid \theta)$ belongs to the Kullback-Leibler support of the prior. They are all weak requirements, satisfied by the most common kernels. Next Lemma shows that they are satisfied under assumptions $B1-B3$.
\begin{Lem}\label{th:WeakConsDens}
Suppose the kernel $k(x \mid \theta)$ satisfies assumptions $B1-B3$. Then $H1-H4$ are fulfilled.
\end{Lem}

\begin{proof}
Assumption $H3$ can be easily deduced from $B1$ and \eqref{data_generating_general}. As regards $H1$, since $\sup_{\theta \in \Theta, x \in \mathbb{X}}k(x \mid \theta) < \infty$, we have $\mathbb{B} = \emptyset$. Moreover, fix $\bar{\theta} \in \mathbb{R}$. If $x > \theta +b$, choose
\[
A(\bar{\theta}, x) = \left(\bar{\theta}-\frac{x-\bar{\theta}-b}{2}, \bar{\theta}+\frac{x-\bar{\theta}-b}{2} \right),
\]
so that $x > \theta+b$ that implies $ k(x \mid \theta) = 0$ for every $\theta \in A(\bar{\theta}, x)$. Similarly, if $x < \theta +a$, choose
\[
A(\bar{\theta}, x) = \left(\bar{\theta}-\frac{\bar{\theta}+a-x}{2}, \bar{\theta}+\frac{\bar{\theta}+a-x}{2} \right).
\]
Finally, if $x \in (\bar{\theta}+a, \bar{\theta}+b)$, denoting $d = \min \{\bar{\theta}+b-x, x-\bar{\theta}-a\}$, choose
\[
A(\bar{\theta}, x) = \left(\bar{\theta}-\frac{d}{2}, \bar{\theta}+\frac{d}{2} \right).
\]
Then $k(x \mid \theta) = g(x-\theta)$ for every $\theta \in A(\bar{\theta}, x)$ and $g$ is continuous on $(a,b)$, by $B2$. Thus we can find the required neighborhood $A(\bar{\theta}, x)$ for every $x \not \in \{\bar{\theta}+a, \bar{\theta}+b\}$, that is for $\mu$-almost every $x$, since $\mu$ is the Lebesgue measure. Therefore $H1$ is satisfied.

$H2$ follows since $\theta$ is a location parameter and $\bar{\Theta} = \Theta$. We are left to show that $H4$ is satisfied: we prove the case $t = 1$ and the general setting follows similarly. 

Recall that assumptions $B1-B3$ can be rewritten as $T1-T4$ in the proof of Theorem \ref{th:ConsBounded} and let $f(x) = k(x\mid \theta^*)$ be the density function of $P$. Fix $\delta > 0$, $\epsilon > 0$ and denote $r= 1-exp(\epsilon/4)$. Define the set
\begin{equation}\label{def_F}
\begin{aligned}
\mathbb{F}(\delta, r) := \biggl\{p(x)=\sum_{j=1}^{\infty} q_j k(x \mid \theta_j) \, : \, &q_1 \in [1-r, 1], q_2 \in [r/2, 1],\\
& 0 \leq \theta^\ast -\theta_1 \leq \delta,0 \leq \theta_2 -\theta^\ast \leq \delta \biggr\}.
\end{aligned}
\end{equation}
We denote $[a_j, b_j] := [a+\theta_j, b+\theta_j]$, with $j \geq 1$, and similarly $[a^*, b^*] := [a + \theta^*, b+\theta^*]$. Then we can choose $\delta$ small enough such that
\[
[a_1,b_1] \cup [a_2, b_2] \supseteq [a^\ast, b^\ast],
\]
for every $\theta_1$ and $\theta_2$ as in \eqref{def_F}. Moreover, for every $x \in S_1 := [a_1, b_1] \cap [a^\ast, b^\ast]$ we have
\[
\begin{aligned}
\log\left\{\frac{g(x-\theta^\ast)}{q_1 g(x-\theta_1)} \right\}&=-\log(q_1)+ \log\left\{\frac{g(x-\theta^\ast)}{g(x-\theta_1)}\right\} \leq \epsilon/4 + \log\left\{\frac{g(x-\theta^\ast)}{g(x-\theta_1)}\right\}\\
& \leq \epsilon/4+R|\theta^*-\theta_1|\,
\end{aligned}
\]
with $R > 0$ as in $T2$. Therefore we can choose $\delta$ small enough so that
\begin{equation}\label{second_delta}
\log\left\{\frac{g(x-\theta^\ast)}{q_1 g(x-\theta_1)} \right\} < \frac{\epsilon}{2}
\end{equation}
for every $x \in S_1$. Similarly, we can choose $\delta$ small enough so that for every $x \in S_2 :=[a^\ast, b^\ast] \setminus [a_1, b_1]$ we have
\begin{equation}\label{third_delta}
\int_{S_2 } g(x-\theta^\ast) \log\bigg\{\frac{g(x-\theta^\ast)}{q_{2} g(x-\theta_2)} \bigg\} \mathrm{d}x  < \frac{\epsilon}{2}.
\end{equation}
Indeed, since $g(x-\theta^*) \leq M$ and $m \leq g(x-\theta_2)$ for every $x$ in $S_2$, with $m$ and $M$ as in $T1$, we have
\[
g(x-\theta^\ast) \log\bigg\{\frac{g(x-\theta^\ast)}{q_2 g(x-\theta_2)} \bigg\} < M \log\{2M/(mr) \},
\]
and $S_2$ has arbitrarily small length with $\delta$ small enough. For every $p \in \mathbb{F}(\delta, r)$, by applying \eqref{second_delta} and \eqref{third_delta}, we have
\begin{align*}
&\int_{a^*}^{b^*} g(x-\theta^\ast) \log\bigg\{\frac{g(x-\theta^\ast)}{\sum_{j=1}^{\infty} q_{j} g(x-\theta_{j})} \bigg\}  \mathrm{d}x =\\
&\int_{S_1} g(x-\theta^\ast) \log\bigg\{\frac{g(x-\theta^\ast)}{\sum_{j=1}^{\infty} q_{j} g(x-\theta_{j})} \bigg\}  \mathrm{d}x +
\int_{S_2} g(x-\theta^\ast) \log\bigg\{\frac{g(x-\theta^\ast)}{\sum_{j=1}^{\infty} q_{j} g(x-\theta_{j})} \bigg\}  \mathrm{d}x \le\\
& \int_{S_1} g(x-\theta^\ast) \log\bigg\{\frac{g(x-\theta^\ast)}{ q_{1} g(x-\theta_{1})} \bigg\}  \mathrm{d}x +
\int_{S_2} g(x-\theta^\ast) \log\bigg\{\frac{g(x-\theta^\ast)}{ q_{2} g(x-\theta_{2})} \bigg\}  \mathrm{d}x \le \epsilon.
\end{align*}
Thus, $\mathbb{F}(\delta, r)\subseteq K_\epsilon(P)$ for $\delta$ small enough.
Moreover, since $\theta^*$ belongs to the support of $Q_0$ and the Dirichlet process prior has full weak support on the space of probability weights $\{ q_j \}_j$, we have that
\[
\text{pr}\{h \in K_\epsilon(P)\} \geq \text{pr}\{h \in \mathbb{F}(\delta, r)\} >0,
\]
as desired.
\end{proof}
The proof of Proposition \ref{undersmoothing} will rely on the following Lemma.
\begin{Lem}\label{lemma_weak_convergence}
Let assumption $H4$ be satisfied and let $P \in \mathbb{G}_t$ with parameters $\theta_1^*, \dots, \theta_t^*$ belonging to the support of $Q_0$. Assume there exists $\mathcal{U}$ weak neighborhood of $P$ such that $\mathcal{U} \cap \mathbb{G}_s = \emptyset$ for every $s < t$. Then
\[
\text{pr}\left(K_n < t \mid X_{1:n} \right) \to 0,
\]
in $P^{(\infty)}$-probability as $n \to \infty$.
\end{Lem}
\begin{proof}
By assumption $H4$, the posterior distribution is consistent at $P$ under the weak topology, in virtue of Schwartz theorem (see e.g.Theorem $6.16$ and Example $6.20$ in \cite{Subhashis2017}), so that
\begin{equation}\label{limit}
\text{pr}(h \in \mathcal{U}^c \mid X_{1:n}) \to 0,
\end{equation}
in $P^{(\infty)}$-probability as $n \to \infty$. Moreover, we have
\[
\text{pr}(h \in \mathcal{U}^c \mid X_{1:n}) \geq \text{pr}(h \in \mathcal{U}^c \mid X_{1:n}, K_n < t)\text{pr}\left(K_n < t \mid X_{1:n} \right).
\]
Notice that, conditional on $K_n < t$, the domain of the posterior distribution is a subset of $\cup_{s < t}\mathbb{G}_s$. Thus we have $\text{pr}(h \in \mathcal{U}^c \mid X_{1:n}, K_n < t) = 1$ and
\[
\text{pr}(h \in \mathcal{U}^c \mid X_{1:n}) \geq \text{pr}\left(K_n < t \mid X_{1:n} \right).
\]
The result follows from \eqref{limit}.
\end{proof}
We need two technical Lemmas.

\begin{Lem}\label{first_auxiliary}
Assume a sequence $\{f_i\}_{i = 1}^\infty \subset \cup_{s < t}\mathbb{G}_s$ is such that $f_i \to f \in \mathcal{P}(\mathbb{X})$ weakly as $i \to \infty$. Then there exist $s' < t$ and a sequence $\{ f'_i\}_{i = 1}^\infty \subset \mathbb{G}_{s'}$ such that $f'_i \to f$ weakly as $i \to \infty$.
\end{Lem}
\begin{proof}
Define
\[
a_s := \sup \{i \geq 1 \, : \, f_i \in \mathbb{G}_s  \}
\]
with $ s < t$. By construction, there exists $s'$ such that $a_{s'} = \infty$ and $\{ f'_i\}$ is the subsequence of elements of $\{ f_i\}$ that belong to $\mathbb{G}_{s'}$.
\end{proof}
\begin{Lem}\label{second_auxiliary}
Let $\left\{ f_i = \sum_{j= 1}^s q_{j,i}k(\cdot \mid \theta_{j,i})\right\}_{i = 1}^\infty \subset \mathbb{G}_s$ be such that $f_i \to f \in \mathcal{P}(\mathbb{X})$ weakly as $i \to \infty$. Then there exist $s' \leq s$ and a sequence $\{ f'_i\}_{i = 1}^\infty \subset \mathbb{G}_{s'}$ such that $f'_i \to f$ weakly as \\$i \to \infty$ and
\[
\lim \inf_i q'_{j,i} > 0
\]
for every $j = 1, \dots, s'$.
\end{Lem}
\begin{proof}
If $\lim \inf_i q_{j,i} = 0$ for every $j = 1, \dots, s$, the statement is true by taking $s := s'$ and $f'_i := f_i$ for every $i \geq 1$. Then assume there exists $l$ such that $\lim \inf_i q_{l,i} = 0$. Consider a subsequence $\{ \tilde{f}_i\}_{i = 1}^\infty$, with weights $\{\tilde{q}_{j,i}\}_i$ and parameters $\{\tilde{\theta}_{j,i}\}_i$, such that $\lim_i \tilde{q}_{l,i} = 0$ and define
\[
f'_i(x) = \sum_{j \neq l} \frac{\tilde{q}_{j,i}}{\sum_{r \neq l}\tilde{q}_{r,i}}k(x \mid \tilde{\theta}_{j,i}),
\]
where $\sum_{r \neq l}\tilde{q}_{r,i} \to 1$, by construction. Let $A \subset \mathbb{X}$, then
\[
\begin{aligned}
\left \lvert \int_A\tilde{f}_i(x) \mu(\d x) - \int_Af'_i(x) \mu(\d x) \right \rvert &=  \sum_{j \neq l}\left(\frac{\tilde{q}_{j,i}}{\sum_{r \neq l}\tilde{q}_{r,i}}- \tilde{q}_{j,i}\right)\int_Ak(x \mid \tilde{\theta}_{j,i}) \mu(\d x) \\
&+\tilde{q}_{l,i}\int_Ak(x \mid \tilde{\theta}_{l,i}) \mu(\d x)\leq \sum_{j \neq l}\left(\frac{\tilde{q}_{j,i}}{\sum_{r \neq l}\tilde{q}_{r,i}}- \tilde{q}_{j,i}\right)+\tilde{q}_{l,i} \to 0,
\end{aligned}
\]
as $i \to \infty$. Therefore, since $A$ is arbitrary and $\{ \tilde{f}_i\}$ converges to $f$, also  $\{ f'_i\}$ converges weakly to $f$ and $\{ f'_i\}_{i = 1}^\infty \in \mathbb{G}_{s-1}$. The result follows by applying recursively the above procedure for every $l$ satisfying $\lim \inf_i q_{l,i} = 0$.
\end{proof}
\begin{proof}[Proof of Proposition \ref{undersmoothing}]% (under the assumptions $H1$--$H3$)]
By Lemma \ref{th:WeakConsDens} we can assume $H1-H4$ and by Lemma \ref{lemma_weak_convergence}, it suffices to prove the existence of a weak neighborhood $\mathcal{U}$ of $P$ such that $\mathcal{U} \cap \mathbb{G}_s = \emptyset$, for every $s < t$. Assume by contradiction that no such $\mathcal{U}$ exists. Then, there exists a sequence $\{f_i\} \in \cap_{s <t}\mathbb{G}_s$ such that $f_i \to f$ weakly, as $i \to \infty$, where $f$ is the density of $P$. By Lemmas \ref{first_auxiliary} and \ref{second_auxiliary} we can assume without loss of generality that $\{f_i\} \in \mathbb{G}_s$, with $s < t$, and $\lim \inf_i q_{j,i} > 0$ for every $j = 1, \dots, s$. We will consider three scenarios, of which at least one must hold: (i) there exists $l \in \{1, \dots, s\}$ such that $\lim \sup_i ||\theta_{l,i}|| = \infty$, (ii) the sequences $\{\theta_{j,i}\}_{i = 1}^\infty$, with $j = 1, \dots, s$, belong to a compact set $C \subset \Theta \backslash \mathbb{B}$ for $i$ large enough, (iii) the sequences $\{\theta_{j,i}\}_{i = 1}^\infty$, with $j = 1, \dots, s$, belong to a compact set $C \subset \Theta$ and there exists $l \in \{1, \dots, s\}$ such that $ \lim \inf _i\inf_{\theta \in \mathbb{B}} ||\theta_{l, i}-\theta || = 0$.

First consider case (i) and assume there exists $1 \leq l \leq s$ such that $|| \theta_{l,r(i)}|| \to \infty$ as $i \to \infty$ for a suitable subsequence $r(i)$. Fix $0 <\epsilon < \lim \inf_i q_{l,i}$ and choose $K \subset \mathbb{X}$ compact set such that $P(K) > 1-\epsilon/4$. By assumption $H2$ we have
\[
\int_{K^c}f_{r(i)}(x) \mu(\d x) > q_{l,r(i)}\int_{K^c}k(x \mid \theta_{l,r(i)}) \mu(\d x) > \frac{\epsilon}{2},
\]
for $i$ large enough, which contradicts the weak convergence of $\{f_i\}_{i = 1}^\infty$ to $f$.

Second, assume to be in case (ii) and there exists a compact set $C \subset \Theta \backslash \mathbb{B}$ such that $\theta_{i,j} \in C$ for every $i \geq 1$ and $j = 1, \dots, s$. Define the set
\[
\mathbb{D}_s : = \left\{\nu(\d \theta) = \sum_{j = 1}^sq_j\delta_{\theta_j}(\d \theta) \, : \, \theta_j \in C, q_j > 0, \sum_{j = 1}^sq_j = 1  \right\} \subset \mathcal{P}(\Theta).
\]
Since $C$ is compact, we have that $\mathbb{D}_s$ is tight. By Prokhorov's Theorem $\mathbb{D}_s$ is also relatively compact, so that there exists a subsequence $r(i)$ such that 
\[
\nu_{r(i)} = \sum_{j = 1}^sq_{j,r(i)}\delta_{\theta_{j,r(i)}} \to \nu \in \mathcal{P}(\Theta)
\]
weakly as $i \to \infty$. By Lemma $4.1$ in \cite{Cai2020} we have $\nu \in \mathbb{D}_{s}$, so that $\nu = \sum_{j = 1}^s\tilde{q}_j\delta_{\tilde{\theta}_j}$ for some $\tilde{q}_j \in (0,1)$, $\sum_{j = 1}^s\tilde{q}_j = 1$ and $\tilde{\theta}_j \in C$, for $j = 1, \dots, s$. By $H1$ and $C \subset \Theta \backslash \mathbb{B}$, for $\mu$-almost every $x \in \mathbb{X}$, we can find $C_j := C_j(x, \tilde{\theta}_j)$, with $j = 1, \dots, s$, closed neighborhood of $\tilde{\theta}_j$, so that $k(x \mid \theta)$ is continuous as a function of $\theta$, with $\theta \in C_j$. Define $D := \left\{\bigcup_{j = 1}^s C_j\right\}\cap C$ compact set: notice that $D \neq \emptyset$, since $\tilde{\theta}_j \in C \cap C_j$, with $j = 1, \dots, s$. Moreover, by construction, the mapping $\theta \in D \to k(x \mid \theta)$ is continuous and therefore bounded, since $D$ is compact. Since $\nu_i \to \nu$ weakly, as $i \to \infty$, there exists $I$ such that for every $i \geq I$ we have $\theta_{j, r(i)} \in D$, for every $j = 1, \dots, s$. Thus, by definition of weak convergence we have
\[
\sum_{j = 1}^sq_{j,r(i)}k(x \mid \theta_{j, r(i)}) = \int k(x \mid \theta) \nu_{r(i)}(\d \theta) \to \int k(x \mid \theta) \nu(\d \theta) = \sum_{j = 1}^s\tilde{q}_{j}k(x \mid \tilde{\theta}_{j}),
\]
as $i \to \infty$. Since almost sure pointwise convergence of densities implies weak convergence, we have
\[
f_{r(i)} \to \tilde{f} = \sum_{j = 1}^s\tilde{q}_{j}k(\cdot \mid \tilde{\theta}_{j})
\]
weakly as $i \to \infty$. By uniqueness of the weak limit, $\tilde{f}(x) = f(x)$ for $\mu$-almost every $x$, that contradicts $H3$.

Third, consider case (iii). Since $\theta_{j, i} \in C \subset \Theta$ compact set, for every $j = 1, \dots, s$ and $i \geq 1$, there exists a suitable subsequence $r(i)$ such that $\theta_{l, r(i)} \to \bar{\theta}$. Since $\mathbb{B}$ is closed by $H1$, we have that $\bar{\theta} \in \mathbb{B}$. By definition of $\mathbb{B}$, this is not possible if $\mu$ is the counting measure, since $k(x \mid \theta) \leq 1$, for every $x \in \mathbb{X}$ and $\theta \in \Theta$. Thus, let $\mu$ be the Lebesgue measure. Then we can fix $\epsilon > 0$ such that
\[
P(B_{x^*}(\epsilon)) < \frac{\lim \inf_i q_{l,i}}{4},
\]
with $x^*$ as in $H2$. Then by $H2$ we have
\[
\int_{B_{x^*}(\epsilon)}f_{r(i)}(x) \mu(\d x) > q_{l,r(i)}\int_{B_{x^*}(\epsilon)}k(x \mid \theta_{l,r(i)}) \mu(\d x) > \frac{\lim \inf_i q_{l,i}}{2},
\]
for $i $ large enough, that again contradicts the weak convergence of $\{f_i\}_{i = 1}^\infty$ to $f$.
\end{proof}


\section{Proof of Theorem \ref{th:ConsUnif}}
The marginal distribution is available and given by the following lemma.
\begin{Lem}\label{marginal_uniform}
Consider $k$ and $q_0$ 
as in \eqref{unif_framework}. Then it holds
\begin{align*}
m(x_{1:n})  &= \frac{2c - \{\max(x_{1:n},\theta^\ast)-\min(x_{1:n},\theta^\ast)\}}{(2c)^{n+1}},
&(x_{1:n}\in[\theta^\ast-c, \theta^\ast+c]^n).
\end{align*}
\end{Lem}

%\subsection{Proof of Lemma \ref{marginal_uniform}}
\begin{proof}%[Proof of Lemma \ref{marginal_uniform}]
Note that $ x_i\in(\theta-c,\theta+c)$ for all $i \in \{1, \dots, n\}$ if and only if $\theta\in(\max(x_{1:n})-c ,\min(x_{1:n})+c)$. Thus
\begin{align*}
m(x_{1:n}) &= \frac{1}{(2c)^{n+1}} \int_{\Theta}\prod_{i = 1}^n \mathbbm{1}_{(\theta-c,\theta+c)}(x_i) \mathbbm{1}_{(\theta^\ast-c,\theta^\ast+c)}(\theta) \mathrm{d}\theta
\\ &=
\frac{1}{(2c)^{n+1}} \int_{\Theta}\mathbbm{1}_{(\max(x_{1:n})-c,\min(x_{1:n})+c)}(\theta) \mathbbm{1}_{(\theta^\ast-c,\theta^\ast+c)}(\theta) \mathrm{d}\theta
\\&=
\frac{2c- \{\max(x_{1:n},\theta^\ast)-\min(x_{1:n},\theta^\ast)\}}{(2c)^{n+1}}\,.
\end{align*}
%    $P^{(\infty)}-$almost surely.
\end{proof}
Define $\text{Range}(X_A) = \max_{i \in A} \left( X_{i}\right)-\min_{i \in A}\left(X_{i} \right)$. Lemma \ref{marginal_uniform} has an important corollary, that is stated after a technical lemma.




\begin{Lem}\label{lemma_unif2}
Let $A \subset \{1, \dots, n\}$ such that $|A| = a$, Then it holds:
\[
\frac{2c- \{\max(X_{A},\theta^\ast)-\min(X_{A},\theta^\ast)\}}{(2c)^{a+1}} \le  \frac{2c-\text{\rm Range}(X_{A})}{(2c)^{a+1}}.
\]
\end{Lem}
\begin{proof}
The result follows immediately from $\max(X_{A},\theta^\ast) \geq \max(X_{A})$ and $\min(X_{A},\theta^\ast) \leq \min(X_{A})$.
\end{proof}
\begin{Cor}\label{main_cor}
In the setting of \eqref{eq:MixDPM} with $(f, k, q_0)$ as in \eqref{unif_framework}, define\\
 $\Omega_n = \left\{ x \in X^{\infty} \, : \,   \max (x_{1:n}) \geq  \theta^\ast \text{ and } \min (x_{1:n} ) \leq \theta^\ast \right\}$. Then
\begin{align}\label{eq:m_ratio_Unif_s+1}
\frac{\prod_{j=1}^{s+1} m(X_{A_j})}{m(X_{1:n})} \mathbbm{1}_{\Omega_n}(X_{1:\infty}) \le \frac{\prod_{j=1}^{s+1} \{2c-\text{\rm Range}(X_{A_j})\}}{(2c)^s \{ 2c-\text{\rm Range}(X_{1:n})\}},
\end{align}
for every $A \in \tau_{s+1}(n)$ .
\end{Cor}
\begin{proof}
As regards the numerator, apply firstly Lemma \ref{marginal_uniform} and then Lemma \ref{lemma_unif2} to get
\[
m(X_{A_j}) = \frac{2c- \{\max(X_{A_j},\theta^\ast)-\min(X_{A_j},\theta^\ast)\}}{(2c)^{a_j+1}} \leq \frac{2c-\text{Range}(X_{A_j})}{(2c)^{a_j+1}}, \quad j = 1, \dots, s+1\,.
\]
Apply Lemma \ref{marginal_uniform} to $m(x_{1:n})$ for every $x \in \Omega_n$, to get
\[
\begin{aligned}
m(X_{1:n})\mathbbm{1}_{\Omega_n}(X_{1:\infty}) &=\frac{2c- \{\max(X_{1:n},\theta^\ast)-\min(X_{1:n},\theta^\ast)\}}{(2c)^{n+1}}\mathbbm{1}_{\Omega_n}(X_{1:\infty}) \\
&= \frac{2c- \{\max(X_{1:n})-\min(X_{1:n})\}}{(2c)^{n+1}}\mathbbm{1}_{\Omega_n}(X_{1:\infty}), 
\end{aligned}
\]
as desired.
\end{proof}
The lemma below shows that, in order to prove Theorem \ref{th:ConsUnif}, it is sufficient to show $\mathbbm{1}_{\Omega_n}(X_{1:\infty})\sum_{s = 1}^{n-1}\frac{\text{pr}(K_n = s+1 | X_{1:n})}{\text{pr}(K_n = 1 | X_{1:n})} \to 0$ in $P^{(\infty)}$-probability.
\begin{Lem}\label{lemma_unif_tech}
Consider $f$ as in \eqref{unif_framework} and define $\Omega_n = \left\{ x \in X^{\infty} \, : \,   \max (x_{1:n}) \geq  \theta^\ast \text{ and } \min (x_{1:n} ) \leq \theta^\ast \right\}$. Let $\left\{Y_n\right\}$ be a sequence of positive random variables. Thus, $Y_n\mathbbm{1}_{\Omega_n}(X_{1:\infty}) \to 0$ in $P^{(\infty)}$-probability implies $Y_n \to 0$ in $P^{(\infty)}$-probability.
\end{Lem}
\begin{proof}
First of all, by definition of $f$ we have
\[
\max (X_{1:n}) \to  \theta^\ast + c, \quad \min (X_{1:n} ) \to \theta^\ast - c
\]
almost surely with respect to $P^{(\infty)}$ as $n \to \infty$. Then $P^{(\infty)}(\Omega_n) \to 1$, as $n \to \infty$, by definition of $\Omega_n$. Thus, fix $\epsilon > 0$ and notice that
\[
P^{(\infty)}\left(Y_n > \epsilon \right) = P^{(\infty)}\left\{\left(Y_n > \epsilon\right) \cap \Omega_n \right\} + P^{(\infty)}\left\{\left(Y_n > \epsilon\right) \cap \Omega_n^c \right\}.
\]
The first term on the right-hand side goes to $0$, since $Y_n\mathbbm{1}_{\Omega_n}(X_{1:\infty}) \to 0$ in $P^{(\infty)}$-probability, while the second vanishes because $P^{(\infty)}(\Omega_n^c) \to 0$, both as $n \to \infty$.
\end{proof}
Combining Corollary \ref{main_cor} and Lemma \ref{lemma_unif_tech} we are ready to prove Theorem \ref{th:ConsUnif}.

\begin{proof}[Proof of Theorem \ref{th:ConsUnif}]
For every $s \geq 1$ and $A \in \tau_s(n)$, from Corollary \ref{main_cor} we have
\[
\frac{\prod_{j=1}^{s} m(X_{A_j})}{m(X_{1:n})} \mathbbm{1}_{\Omega_n}(X_{1:\infty}) \le \frac{\prod_{j=1}^{s} \{2c-\text{Range}(X_{A_j})\}}{(2c)^{s-1} \{ 2c-\text{Range}(X_{1:n})\}}.
\]
Note that $\{2c-\text{Range}(X_{A_j})\}/(2c) \sim \text{Beta}(2,a_j-1)$ independently for $j=1,\ldots,s$. %and $M=1/\{1-\text{Range}(X_{A_j})\}$, where $1/M \sim \text{Be}(2,n-1)$.
Moreover, recall that if $Z \sim \text{Beta}(\alpha,\beta)$ then for $p > - \alpha$
\begin{equation*}
E(Z^p) = \frac{\Gamma(\alpha+p) \Gamma(\alpha+\beta)}{\Gamma(\alpha+p+\beta) \Gamma(\alpha)}.
\end{equation*}
Thus, by H\"{o}lder's inequality with exponents $3$ and $3/2$ we get
\begin{align*}
E\bigg\{\frac{\prod_{j=1}^{s} m(X_{A_j})}{m(X_{1:n})} \bigg\} & \leq E\left\{\prod_{j=1}^{s} m(X_{A_j})^3 \right\}^{1/3}E\left\{ m(X_{1:n}) ^{-3/2}\right\}^{2/3} \\
&= \bigg\{\frac{\Gamma(5)}{\Gamma(2)}\bigg\}^{s/3} \bigg\{\frac{\Gamma(1/2)}{\Gamma(2)}\bigg\}^{2/3} \bigg\{\prod_{j=1}^{s} \frac{\Gamma(1+a_j)}{\Gamma(a_j+4)}\bigg\}^{1/3} \bigg\{\frac{\Gamma(1+n)}{\Gamma (n-1/2 )}\bigg\}^{2/3}.
\end{align*}
By the recursive definition of the Gamma function and recalling that $\Gamma(1/2)=\pi^{1/2}$, the upper bound above becomes
\begin{align*}
E\bigg\{\frac{\prod_{j=1}^{s} m(X_{A_j})}{m(X_{1:n})} \bigg\} &\leq	24^{s/3} \pi^{1/3} \bigg\{\prod_{j=1}^{s} \frac{\Gamma(1+a_j)}{\Gamma(a_j+4)}\bigg\}^{1/3} \bigg\{\frac{\Gamma(1+n)}{\Gamma(n-1/2)}\bigg\}^{2/3}\\
&= 24^{s/3} \pi^{1/3} \bigg\{\prod_{j=1}^{s} \frac{1}{(a_j+3)(a_j+2)(a_j+1)}\bigg\}^{1/3} \bigg\{ \frac{(n-1/2)\Gamma(1+n)}{\Gamma(n+1/2)} \bigg\}^{2/3}.
\end{align*}
%The previous equalities follows by the recursive definition of Gamma function and recalling that $\Gamma(1/2)=\pi^{1/2}$. 
Moreover, exploiting again the recursive definition of the Gamma function, Gautschi's Inequality, i.e.\ $\frac{\Gamma(1+n)}{\Gamma(n+1/2)} \le (n+1)^{1/2}$, and 
%the simple inequalities $n-1/2<n+1/2$ and 
$(n+1)/(a_j+1)<n/a_j$, we have
\begin{align*}
E\bigg\{\frac{\prod_{j=1}^{s} m(X_{A_j})}{m(X_{1:n})} \bigg\} \le 24^{s/3} K \bigg\{\prod_{j=1}^{s} \frac{(n+1)^3}{(a_j+1)^3}\bigg\}^{1/3}
\le 24^{s/3} K  \bigg(\frac{n^3}{\prod_{j=1}^{s} a_i^3} \bigg)^{1/3} = 24^{s/3} K  \frac{n}{\prod_{j=1}^{s} a_j}.
\end{align*}
Thus, applying Lemma \ref{suff_expectation} and Lemma \ref{induction} with $p=2$ and $C=4 \zeta(2) < 7$ we get
\begin{equation*}
E\{R(n,1,s)\} \le \frac{24^{s/3} K}{s!} \sum_{\bm{a} \in \mathcal{F}_s(n)} \bigg(\frac{n}{\prod_{j=1}^s a_j} \bigg)^2 < \frac{C^{s-1} 24^{s/3} K}{s!},
\end{equation*}
where $R(n,1,s)$ is defined as in \eqref{newratios}. From Corollary \ref{cor_C} we have
\[
C(n,1,s+1) \leq \frac{G\Gamma(2+\beta)2^ss}{\epsilon}E(\alpha^s)\log\{n/(1+\epsilon)\}^{-1}, \quad n \geq 4\,.
\]
Thus, combining the inequalities above with \eqref{newratios} and assumption $A3$ we have
\[
\begin{aligned}
E \biggl\{\mathbbm{1}_{\Omega_n}(X_{1:\infty})&\sum_{s = 1}^{n-1}\frac{\text{pr}(K_n = s+1 | X_{1:n})}{\text{pr}(K_n = 1 | X_{1:n})}\biggr\} = \sum_{s = 1}^{n-1}C(n,1,s+1)E\{\mathbbm{1}_{\Omega_n}(X_{1:\infty})R(n,1,s+1)\}\\
&  \leq \frac{24^{1/3}DGK\Gamma(2+\beta)}{\epsilon \log\{n/(1+\epsilon)\}}\underbrace{\sum_{ s = 1}^{n-1}\frac{s (2C24^{1/3})^{s}\rho^{-s}\Gamma(\nu +s+1)}{(s+1)!}}_{< \infty} \to 0\qquad\hbox{as }n\to\infty\,,
\end{aligned}
\]
where finiteness follows from  $\rho \ge 38 > 24^{1/3} \times 2C$. This implies that
\begin{equation*}
\sum_{s = 1}^{n-1}\frac{\text{pr}(K_n = s+1 | X_{1:n})}{\text{pr}(K_n = 1 | X_{1:n})} \rightarrow 0
\end{equation*}
in $L^1$ and thus in $P^{(\infty)}$-probability as $n\to\infty$. Lemma \ref{lemma_unif_tech} with $Y_n = \sum_{s = 1}^{n-1}\frac{\text{pr}(K_n = s+1 | X_{1:n})}{\text{pr}(K_n = 1 | X_{1:n})}$ concludes the proof.
\end{proof}


\section{Proof of Theorem \ref{th:ConsNormDelta}}% (Gaussian mixtures)}
We first need the following result.
\begin{Lem}\label{lemma:m_mratio_Gauss}
Let $k$ and $q_0$ be as in \eqref{eq:NormDelta} and $x_1=\cdots=x_n=\theta^\ast$ for some $\theta^\ast \in\R$. Then
\begin{align*}
\frac{\prod_{j=1}^{s}m(x_{A_j})}{m(x_{1:n})} &= \left\{\frac{n+1}{\prod_{j=1}^{s} (a_j+1)} \right\}^{1/2} \exp\bigg\{\frac{{\theta^\ast}^2}{2} \bigg(- \frac{n^2}{n+1} + \sum_{j=1}^{s} \frac{a_j^2}{a_j+1} \bigg) \bigg\}
< \left(\frac{n}{\prod_{j=1}^{s} a_j} \right)^{1/2}\,,
\end{align*}
for every $s=1,\dots,n$ and every partition $A=\{A_1,\ldots,A_s\} \in \tau_s(n)$.
\end{Lem}
\begin{proof}
Since the marginal likelihood can be rewritten as
\begin{equation*}
m(x_{A_j})= (a_j+1)^{-1/2} q_0(\theta^\ast)^{a_j} \exp\bigg\{\frac{{\theta^\ast}^2}{2} \frac{a_j^2}{a_j+1} \bigg\},
\end{equation*}
the first equality is obtained. The inequality follows from
\[
\begin{aligned}
- \frac{n^2}{n+1} + \sum_{j=1}^{s} \frac{a_j^2}{a_j+1} &= n - \frac{n^2}{n+1} + \sum_{j=1}^{s} \bigg(\frac{a_j^2}{a_j+1} -a_j \bigg) =  \frac{n}{n+1} - \sum_{j=1}^{s} \frac{a_j}{a_j+1} = \\
& = \sum_{j=1}^{s} a_j \bigg(\frac{1}{n+1} -\frac{1}{a_j+1} \bigg) \le 0
\end{aligned}
\]
and
\[
\frac{n+1}{\prod_{j=1}^{s} (a_j+1)} \leq \frac{n}{\prod_{j=1}^{s} a_j},
\]
which easily follows from $a_j \leq n$, for every $j = 1, \dots, s$.
\end{proof}

\begin{proof}[Proof of Theorem \ref{th:ConsNormDelta}]
First, we study $R(n,1,s)$ as defined in \eqref{newratios}. Since all the observations are almost surely equal, we have
\begin{equation*}
R(n,1,s) = \sum_{\bm{a} \in \mathcal{F}_s(n)} \frac{n}{s! \prod_{j=1}^{s} a_j} \frac{\prod_{j=1}^s m(X_{A^{\bm{a}}_j})}{m(X_{1:n})}\,,
\end{equation*}
where $A^{\bm{a}}$ is an arbitrary partition in $\tau_s(n)$ such that $|A_j^{\bm{a}}| = a_j$ for $j=1,\dots,s$. By application of Lemma \ref{lemma:m_mratio_Gauss} and Lemma \ref{induction} with $p=3/2$, it turns out that the constant $C = 2^{\frac{3}{2}}\zeta\left(\frac{3}{2} \right) < 8$ is such that
\begin{equation*}
R(n,1,s) < \frac{1}{s!} \sum_{\bm{a} \in \mathcal{F}_s(n)} \left(\frac{n}{\prod_{j=1}^{s} a_j} \right)^{3/2} < \frac{C^{s-1}}{s!}.
\end{equation*}
From Corollary \ref{cor_C} we have
\begin{equation}\label{eq:thm1_1}
C(n,1,s+1) \leq \frac{G\Gamma(2+\beta)2^ss}{\epsilon}E(\alpha^s)\log\{n/(1+\epsilon)\}^{-1}, \quad n \geq 4\,.
\end{equation}
Thus, combining the inequalities above with \eqref{newratios} and assumption $A3$ we have
\begin{equation}\label{eq:thm1_2}
\begin{aligned}
\sum_{s = 1}^{n-1}\frac{\text{pr}(K_n = s+1 | X_{1:n})}{\text{pr}(K_n = 1 | X_{1:n})} &= \sum_{s = 1}^{n-1}C(n,1,s+1)R(n,1,s+1)\\
&  \leq \frac{DG\Gamma(2+\beta)}{\epsilon \log\{n/(1+\epsilon)\}}\underbrace{\sum_{ s = 1}^{n-1}\frac{s (2C)^{s}\rho^{-s}\Gamma(\nu +s+1)}{(s+1)!}}_{< \infty} \to 0\qquad\hbox{as }n\to\infty\,,
\end{aligned}
\end{equation}
where the finiteness follows from $\rho > 16 > 2C$. Then we conclude applying a variation of Lemma \ref{sufficient} with equalities and limits in probability replaced by almost sure equalities and limits (the proof of Lemma \ref{sufficient} extends trivially to that case).
\end{proof}

\section{Proof of Proposition \ref{prop:PostAlpha}}
\begin{proof}
Under \eqref{eq:MixDPM}, for every $\epsilon > 0$ we have
\begin{align*}
\text{pr}(\alpha < \epsilon \mid X_{1:n}) &= \sum_{s=1}^{n} \text{pr}(\alpha < \epsilon \mid  K_n=s) \,\ \text{pr}(K_n=s \mid X_{1:n}) =\\
&\ge \text{pr}(\alpha < \epsilon \mid K_n=t) \,\ \text{pr}(K_n=t \mid X_{1:n}).
\end{align*}
By assumption, $\text{pr}(K_n=t \mid X_{1:n}) \to 1$ in $P^{(\infty)}$-probability as $n \to \infty$. Moreover, by Proposition \ref{prop:Prior} with $s = 1$ we get
\[
E(\alpha \mid K_n=t) = C(n, t, t+1) \rightarrow 0,
\]
as $n \to \infty$. It follows $\text{pr}(\alpha < \epsilon \mid K_n=t) \to 1$ in $P^{(\infty)}$-probability as $n \to \infty$, as desired.
\end{proof}

\section*{References}

Cai, D., Campbell, T., and Broderick, T. (2021). Finite mixture models do not reliably 
\hspace*{0.2cm} learn the number of components. \emph{38th International Conference on Machine Learning} \hspace*{0.2cm} {\bf{139}}, 1158--1169.\\
\\
Ghosal, S.\ and Van Der Vaart, A.\ W. (2017). \emph{Fundamentals of Nonparametric Bayesian \hspace*{0.2cm} Inference}. Cambridge
University Press.


\end{document}
